% $Id$

\documentclass[11pt]{scrartcl}
\usepackage{a4wide,verbatim,graphicx}
\usepackage{pdfsetup}
\usepackage{html}
\usepackage{microtype}

\newcommand{\bl}{\latex{\symbol{123}}\html{\{}}
\newcommand{\br}{\latex{\symbol{125}}\html{\}}}

\newcommand{\xsmall}{}
\newcommand{\xtiny}{\small}

\newcommand{\trit}{\em}

\newcommand{\ver}{1.5}
\newcommand{\nextver}{1.6}
\newcommand{\latestunicodever}{6.3}

\begin{document}
\latex{{\centerline{\hspace*{0mm}\includegraphics[scale=0.6]{logo}}}}

\html{
\begin{rawhtml}
<CENTER>
<A NAME="TOP"></a>
<A HREF="http://www.jflex.de"><IMG SRC="logo.png" BORDER=0 HEIGHT=160 WIDTH=577></a>
</CENTER>
\end{rawhtml}
}

\begin{center}
\sffamily
{\Large The Fast Lexical Analyser Generator}\\
\smallskip\smallskip
Copyright \copyright 1998--2013 by \htmladdnormallink{Gerwin Klein}{http://www.doclsf.de},
Steve Rowe, and R\'egis D\'ecamps.
\\
\bigskip
{\Huge \sffamily \bfseries JFlex User's Manual}\\
\bigskip
Version \ver, {\today}
\end{center}

\tableofcontents
\vfill
\newpage
\parindent0pt\parskip1ex
\pagestyle{headings}
\section{Introduction\label{Intro}}
JFlex is a lexical analyser generator for Java\footnote{Java is a trademark of 
Sun Microsystems, Inc., and refers to Sun's Java programming language. 
JFlex is not sponsored by or affiliated with Sun Microsystems, Inc.}
written in Java. It is also a rewrite of the very useful tool JLex \cite{JLex} which
was developed by Elliot Berk at Princeton University. As Vern Paxson states
for his C/C++ tool flex \cite{flex}: they do not share any code though.

\subsection{Design goals}
The main design goals of JFlex are:
\begin{itemize}
\item {\bf Full unicode support}
\item {\bf Fast generated scanners }
\item {\bf Fast scanner generation}
\item {\bf Convenient specification syntax}
\item {\bf Platform independence}
\item {\bf JLex compatibility}
\end{itemize}

\subsection{About this manual}
This manual gives a brief but complete description of the tool JFlex. It
assumes that you are familiar with the issue of lexical analysis. The references \cite{Aho},
\cite{Appel}, and \cite{Maurer} provide a good introduction to this topic.

The next section of this manual describes \htmlref{\emph{installation procedures}}{Installing} 
for JFlex. If you never worked with JLex or
just want to compare a JLex and a JFlex scanner specification you
should also read \htmlref{\emph{Working with JFlex - an example}}{Example} 
(section \ref{Example}).  All options and the complete
specification syntax are presented in 
\htmlref{\emph{Lexical specifications}}{Specifications} (section \ref{Specifications});
\htmlref{\emph{Encodings, Platforms, and Unicode}}{sec:encodings} (section \ref{sec:encodings})
provides information about scanning text vs.~binary files.
If you are interested in performance
considerations and comparing JLex with JFlex speed, 
\htmlref{\emph{a few words on performance}}{performance} (section \ref{performance})
might be just right for you. Those who want to
use their old JLex specifications may want to check out section \ref{Porting}
\htmlref{\emph{Porting from JLex}}{Porting} to avoid possible problems
with not portable or non standard JLex behaviour that has been fixed in
JFlex. Section \ref{lexport} talks about porting scanners from the
Unix tools lex and flex.  Interfacing JFlex scanners with the LALR
parser generators CUP and BYacc/J is explained in \htmlref{\emph{working
  together}}{WorkingTog} (section \ref{WorkingTog}).  Section \ref{Bugs}
\htmlref{\emph{Bugs}}{Bugs} gives a list of currently known active bugs. 
The manual concludes with notes about 
\htmlref{\emph{Copying and License}}{Copyright} (section \ref{Copyright}) and
\htmlref{references}{References}.


\section{Installing and Running JFlex\label{Installing}}

\subsection{Installing JFlex}

\subsubsection{Windows\label{install:windows}}
To install JFlex on Windows 95/98/NT/XP, follow these three steps:
\begin{enumerate}
\item
Unzip the file you downloaded into the directory you want JFlex in (using
something like 
\htmladdnormallinkfoot{WinZip}{http://www.winzip.com}).
If you unzipped it to say \verb+C:\+, the following directory structure
should be generated:
{\xsmall
\begin{verbatim}
C:\JFlex\                          
      +--bin\                   (start scripts) 
      +--doc\                   (FAQ and manual)
      +--examples\ 
               +--binary\       (scanning binary files)
               +--byaccj\       (calculator example for BYacc/J)
               +--cup\          (calculator example for cup)
               +--interpreter\  (interpreter example for cup)
               +--java\         (Java lexer specification) 
               +--simple\       (example scanner)
               +--simple-maven\ (example with maven)
               +--standalone\   (a simple standalone scanner)
               +--standalone-maven\  (above with maven)
      +--lib\                   (the precompiled classes) 
      +--src\ 
          +--jflex\             (source code of JFlex) 
          +--jflex\gui          (source code of JFlex UI classes)
          +--java_cup\runtime\  (source code of cup runtime classes) 
\end{verbatim}
}

\item
Edit the file {\bf \verb+bin\jflex.bat+}
(in the example it's \verb+C:\JFlex\bin\jflex.bat+)
such that

\begin{itemize}
\item
  {\bf \texttt{JAVA\_HOME}} contains the directory where your Java JDK is installed
  (for instance \verb+C:\java+) and
\item
  {\bf \texttt{JFLEX\_HOME}} the directory that contains JFlex (in the example:
  \verb+C:\JFlex+) 
\end{itemize}
 
\item
Include the \verb+bin\+ directory of JFlex in your path.
(the one that contains the start script, in the example: \verb+C:\JFlex\bin+).
\end{enumerate}

\subsubsection{Unix with tar archive}

To install JFlex on a Unix system, follow these two steps:
\begin{itemize}
\item
  Decompress the archive into a directory of your choice
  with GNU tar, for instance to \texttt{/usr/share}:

  \texttt{tar -C /usr/share -xvzf jflex-\ver.tar.gz}
  
  (The example is for site wide installation. You need to
  be root for that. User installation works exactly the
  same way---just choose a directory where you have write
  permission)

\item
  Make a symbolic link from somewhere in your binary 
  path to \texttt{bin/jflex}, for instance:

  \texttt{ln -s /usr/share/JFlex/bin/jflex /usr/bin/jflex}

  If the Java interpreter is not in your binary path, you
  need to supply its location in the script \texttt{bin/jflex}.
\end{itemize}

You can verify the integrity of the downloaded file with
the MD5 checksum available on the \htmladdnormallink{JFlex download page}{http://www.jflex.de/download.html}. 
If you put the checksum file in the same directory
as the archive, you run:

\verb+md5sum --check +\texttt{jflex-\ver.tar.gz.md5}

It should tell you

\texttt{jflex-\ver.tar.gz: OK}


\subsection{Running JFlex}
You run JFlex with:

\texttt{jflex <options> <inputfiles>}

It is also possible to skip the start script in \verb+bin/+
and include the file \verb+lib/JFlex.jar+
in your \texttt{CLASSPATH} environment variable instead.

Then you run JFlex with:

\texttt{java jflex.Main <options> <inputfiles>}

or with:

\texttt{java -jar \texttt{JFlex.jar} <options> <inputfiles>}

The input files and options are in both cases optional. If you don't
provide a file name on the command line, JFlex will pop up a window
to ask you for one.

JFlex knows about the following options:

\begin{description}
\item \verb+-d <directory>+\\
  writes the generated file to the directory \verb+<directory>+

\item \verb+--skel <file>+\\
  uses external skeleton \verb+<file>+. This is mainly for JFlex
  maintenance and special low level customisations. Use only when you
  know what you are doing! JFlex comes with a skeleton file in the
  \texttt{src} directory that reflects exactly the internal, pre-compiled
  skeleton and can be used with the \texttt{-skel} option.

\item \verb+--nomin+\\
  skip the DFA minimisation step during scanner generation.

\item \verb+--jlex+\\
  tries even harder to comply to JLex interpretation of specs.

\item \verb+--dot+\\
  generate graphviz dot files for the NFA, DFA and minimised 
  DFA. This feature is still in alpha status, and not
  fully implemented yet. 

\item \verb+--dump+\\
  display transition tables of NFA, initial DFA, and minimised DFA 

\item \verb+--legacydot+\\
  dot (\texttt{.}) meta character matches \verb+[^\n]+ instead of\\
  \verb+[^\n\r\u000B\u000C\u0085\u2028\u2029]+
  
\item \verb+--noinputstreamctor+\\
  don't include an InputStream constructor in the generated scanner

\item \verb+--verbose+ or \texttt{-v}\\
  display generation progress messages (enabled by default)

\item \verb+--quiet+ or \texttt{-q}\\
  display error messages only (no chatter about what JFlex is
  currently doing)

\item \verb+--time+\\
  display time statistics about the code generation process
  (not very accurate)

\item \verb+--version+\\
  print version number 

\item \verb+--info+\\
  print system and JDK information (useful if you'd like
  to report a problem)
  
\item \verb+--unicodever <ver>+\\
  print all supported properties for Unicode version \verb+<ver>+

\item \verb+--pack+\\
  use the \%pack code generation method by default

\item \verb+--table+\\
  use the \%table code generation method by default\\
  {\bf Deprecated - will be removed in JFlex \nextver}

\item \verb+--switch+\\
  use the \%switch code generation method by default\\
  {\bf Deprecated - will be removed in JFlex \nextver}

\item \verb+--help+ or \texttt{-h}\\
  print a help message explaining options and usage of JFlex.
\end{description}

\section{A simple Example: How to work with JFlex\label{Example}}
To demonstrate what a lexical specification with JFlex looks like, this
section presents a part of the specification for the Java language.
The example does not describe the whole lexical structure of Java programs,
but only a small and simplified part of it (some keywords, some operators,
comments and only two kinds of literals). It also shows how to interface
with the LALR parser generator CUP \cite{CUP} and therefore
uses a class \texttt{sym} (generated by CUP), where integer constants for
the terminal tokens of the CUP grammar are declared. JFlex comes with a
directory \texttt{examples}, where you can find a small standalone scanner
that doesn't need other tools like CUP to give you a running example.
The "\texttt{examples}" directory also contains a {\trit complete} JFlex 
specification of the lexical structure of Java programs together with the 
CUP parser specification for Java by 
\htmladdnormallink{C. Scott Ananian}{mailto:cananian@alumni.princeton.edu}, obtained
from the CUP \cite{CUP} web site (it was modified to interface with the JFlex scanner). 
Both specifications adhere to the Java Language Specification \cite{LangSpec}.

{\xtiny
\label{CodeTop}\begin{verbatim}
/* JFlex example: part of Java language lexer specification */
import java_cup.runtime.*;

/**
 * This class is a simple example lexer.
 */
%%
\end{verbatim}
\label{CodeOptions}\begin{verbatim}
%class Lexer
%unicode
%cup
%line
%column
\end{verbatim}
\label{CodeScannerCode}\begin{verbatim}
%{
  StringBuffer string = new StringBuffer();

  private Symbol symbol(int type) {
    return new Symbol(type, yyline, yycolumn);
  }
  private Symbol symbol(int type, Object value) {
    return new Symbol(type, yyline, yycolumn, value);
  }
%}
\end{verbatim}
\label{CodeMacros}\begin{verbatim}
LineTerminator = \r|\n|\r\n
InputCharacter = [^\r\n]
WhiteSpace     = {LineTerminator} | [ \t\f]

/* comments */
Comment = {TraditionalComment} | {EndOfLineComment} | {DocumentationComment}

TraditionalComment   = "/*" [^*] ~"*/" | "/*" "*"+ "/"
EndOfLineComment     = "//" {InputCharacter}* {LineTerminator}
DocumentationComment = "/**" {CommentContent} "*"+ "/"
CommentContent       = ( [^*] | \*+ [^/*] )*

Identifier = [:jletter:] [:jletterdigit:]*

DecIntegerLiteral = 0 | [1-9][0-9]*
\end{verbatim}
\label{CodeStateDecl}\begin{verbatim}
%state STRING

%%
\end{verbatim}
\label{CodeRulesYYINITIAL}\begin{verbatim}
/* keywords */
<YYINITIAL> "abstract"           { return symbol(sym.ABSTRACT); }
<YYINITIAL> "boolean"            { return symbol(sym.BOOLEAN); }
<YYINITIAL> "break"              { return symbol(sym.BREAK); }
\end{verbatim}
\label{CodeRulesBunch}\begin{verbatim}
<YYINITIAL> {
  /* identifiers */ 
  {Identifier}                   { return symbol(sym.IDENTIFIER); }
 
  /* literals */
  {DecIntegerLiteral}            { return symbol(sym.INTEGER_LITERAL); }
  \"                             { string.setLength(0); yybegin(STRING); }

  /* operators */
  "="                            { return symbol(sym.EQ); }
  "=="                           { return symbol(sym.EQEQ); }
  "+"                            { return symbol(sym.PLUS); }

  /* comments */
  {Comment}                      { /* ignore */ }
 
  /* whitespace */
  {WhiteSpace}                   { /* ignore */ }
}
\end{verbatim}
\label{CodeRulesYYtext}\begin{verbatim}
<STRING> {
  \"                             { yybegin(YYINITIAL); 
                                   return symbol(sym.STRING_LITERAL, 
                                   string.toString()); }
  [^\n\r\"\\]+                   { string.append( yytext() ); }
  \\t                            { string.append('\t'); }
  \\n                            { string.append('\n'); }

  \\r                            { string.append('\r'); }
  \\\"                           { string.append('\"'); }
  \\                             { string.append('\\'); }
}
\end{verbatim}
\label{CodeRulesAllStates}\begin{verbatim}
/* error fallback */
[^]                              { throw new Error("Illegal character <"+
                                                    yytext()+">"); }
\end{verbatim}
}

From this specification JFlex generates a \texttt{.java} file with one
class that contains code for the scanner. The class will have a
constructor taking a \texttt{java.io.Reader} from which the input is
read. The class will also have a function \texttt{yylex()} that runs the
scanner and that can be used to get the next token from the input (in this
example the function actually has the name \texttt{next\_token()} because
the specification uses the \texttt{\%cup} switch).

As with JLex, the specification consists of three parts, divided by \texttt{\%\%}:
\begin{itemize}
\item \htmlref{usercode}{ExampleUserCode},
\item \htmlref{options and declarations}{ExampleOptions} and 
\item \htmlref{lexical rules}{ExampleLexRules}.
\end{itemize}

\subsection{Code to include\label{ExampleUserCode}}
Let's take a look at the first section, ``user code'': The text up to the
first line starting with \texttt{\%\%} is copied verbatim to the top 
of the generated lexer class (before the actual class declaration). 
Beside \texttt{package} and \texttt{import} statements there is usually not much 
to do here. If the code ends with a javadoc class comment, the generated class
will get this comment, if not, JFlex will generate one automatically.

\subsection{Options and Macros\label{ExampleOptions}}
The second section ``options and declarations'' is more interesting. It consists
of a set of options, code that is included inside the generated scanner
class, lexical states and macro declarations. Each JFlex option must begin
a line of the specification and starts with a \texttt{\%}. In our example
the following options are used:

\begin{itemize}
\item
  \texttt{\htmlref{\%class Lexer}{CodeOptions}} tells JFlex to give the
  generated class the name ``Lexer'' and to write the code to a file ``\texttt{Lexer.java}''.

\item
  \texttt{\htmlref{\%unicode}{CodeOptions}} defines the set of characters the scanner will 
  work on. For scanning text files, \texttt{\%unicode} should always be used.  The Unicode
  version may be specified, e.g. \texttt{\%unicode 4.1}.  If no version is specified, the
  most recent supported Unicode version will be used - in JFlex \ver, this is
  Unicode \latestunicodever. See also section \ref{sec:encodings} for more information
  on character sets, encodings, and scanning text vs. binary files.

\item
  \texttt{\htmlref{\%cup}{CodeOptions}} switches to CUP compatibility
  mode to interface with a CUP generated parser.

\item
  \texttt{\htmlref{\%line}{CodeOptions}} switches line counting on (the
  current line number can be accessed via the variable \texttt{yyline})

\item
  \texttt{\htmlref{\%column}{CodeOptions}} switches column counting on
  (current column is accessed via \texttt{yycolumn})

\end{itemize}
\label{ExampleScannerCode}

The code included in \texttt{\htmlref{\%\bl  ...\%\br}{CodeScannerCode}}
is copied verbatim into the generated lexer class source.
Here you can declare member variables and functions that are used
inside scanner actions. In our example we declare a \texttt{StringBuffer} ``\texttt{string}''
in which we will store parts of string literals and two helper functions
``\texttt{symbol}'' that create \texttt{java\_cup.runtime.Symbol} objects
with position information of the current token (see section \ref{CUPWork}
\htmlref{\emph{JFlex and CUP}}{CUPWork} 
for how to interface with the parser generator CUP). As JFlex options, both 
\verb+%{+ and \verb+\%}+ must begin a line. 
\label{ExampleMacros}

The specification continues with macro declarations. Macros are
abbreviations for regular expressions, used to make lexical specifications
easier to read and understand. A macro declaration
consists of a macro identifier followed by \texttt{=}, then followed by
the regular expression it represents.  This regular expression may
itself contain macro usages. Although this allows a grammar like specification
style, macros are still just abbreviations and not non terminals -- they
cannot be recursive or mutually recursive. Cycles in macro definitions
are detected and reported at generation time by JFlex.

Here some of the example macros in more detail:
\begin{itemize}
\item
  \texttt{\htmlref{LineTerminator}{CodeMacros}} stands for the regular
  expression that matches an ASCII CR, an ASCII LF or an CR followed by LF.

\item
  \texttt{\htmlref{InputCharacter}{CodeMacros}} stands for all characters
  that are not a CR or LF.

\item
  \texttt{\htmlref{TraditionalComment}{CodeMacros}} is the expression
  that matches the string \texttt{"/*"} followed by a character that
  is not a \texttt{*}, followed by anything that does not contain, but
  ends in \texttt{"*/"}.  As this would not match comments like
  \texttt{/****/}, we add \texttt{"/*"} followed by an arbitrary
  number (at least one) of \texttt{"*"} followed by the closing
  \texttt{"/"}. This is not the only, but one of the simpler
  expressions matching non-nesting Java comments. It is tempting to
  just write something like the expression \texttt{"/*" .* "*/"}, but
  this would match more than we want. It would for instance match the
  whole of \texttt{/* */ x = 0; /* */}, instead of two comments and
  four real tokens. See DocumentationComment and CommentContent for an
  alternative.

\item
  \texttt{\htmlref{CommentContent}{CodeMacros}} matches zero or more
  occurrences of any character except a \texttt{*} or any number of 
  \texttt{*} followed by a character that is not a \texttt{/}

\item
  \texttt{\htmlref{Identifier}{CodeMacros}} matches each string that
  starts with a character of class \texttt{jletter} followed by zero or more characters
  of class \texttt{jletterdigit}. \texttt{jletter} and \texttt{jletterdigit}
  are predefined character classes.  \texttt{jletter} includes all characters for which
  the Java function \texttt{Character.isJavaIdentifierStart} returns \texttt{true} and 
  \texttt{jletterdigit} all characters for that \texttt{Character.isJavaIdentifierPart}
  returns \texttt{true}. 
\end{itemize}
\label{ExampleStateDecl}

The last part of the second section in our
lexical specification is a lexical state declaration: 
\texttt{\htmlref{\%state STRING}{CodeStateDecl}}
declares a lexical state \texttt{STRING} that can be
used in the ``lexical rules'' part of the specification. A state declaration
is a line starting with \texttt{\%state} followed by a space or comma
separated list of state identifiers. There can be more than one line starting
with \texttt{\%state}.

\subsection{Rules and Actions\label{ExampleLexRules}}
The "lexical rules" section of a JFlex specification contains regular expressions
and actions (Java code) that are executed when the scanner matches the
associated regular expression. As the scanner reads its input, it keeps
track of all regular expressions and activates the action of the expression
that has the longest match. Our specification above for instance would with input
"\texttt{breaker}" match the regular expression for \texttt{\htmlref{Identifier}{CodeMacros}}
and not the keyword "\texttt{\htmlref{break}{CodeRulesYYINITIAL}}"
followed by the Identifier "\texttt{er}", because rule \verb+{Identifier}+
matches more of this input at once (i.e. it matches all of it)
than any other rule in the specification. If two regular expressions both 
have the longest match for a certain input, the scanner chooses the action 
of the expression that appears first in the specification. In that way, we 
get for input "\texttt{break}" the keyword "\texttt{break}" and not an 
Identifier "\texttt{break}". 

Additional to regular expression matches, one can use lexical states to
refine a specification. A lexical state acts like a start condition.
If the scanner is in lexical state \texttt{STRING}, only expressions that
are preceded by the start condition \texttt{<STRING>} can be matched.
A start condition of a regular expression can contain more than one lexical
state. It is then matched when the lexer is in any of these lexical states.
The lexical state \texttt{YYINITIAL} is predefined and is also the state
in which the lexer begins scanning. If a regular expression has no start
conditions it is matched in {\trit all} lexical states.
\label{ExampleRulesStateBunch}

Since you often have a bunch of expressions with the same start conditions, 
JFlex allows the same abbreviation as the Unix tool \texttt{flex}:
\begin{verbatim}
<STRING> {
  expr1   { action1 }
  expr2   { action2 }
}
\end{verbatim}
means that both \texttt{expr1} and \texttt{expr2} have start condition \texttt{<STRING>}.
\label{ExampleRulesYYINITIAL}

The first three rules in our example demonstrate the syntax of a regular 
expression preceded by the start condition \texttt{<YYINITIAL>}.

{\tt\htmlref{<YYINITIAL> "abstract"}{CodeRulesYYINITIAL}\verb+  {+ return symbol(sym.ABSTRACT); \verb+}+}

matches the input "\texttt{abstract}" only if the scanner is in its
start state "\texttt{YYINITIAL}". When the string "\texttt{abstract}" is
matched, the scanner function returns the CUP symbol \texttt{sym.ABSTRACT}.
If an action does not return a value, the scanning process is resumed immediately 
after executing the action.
\label{ExampleRulesBunch}

The rules enclosed in

\texttt{\htmlref{<YYINITIAL> \bl \\
\ \ ...\\
\br}{CodeRulesBunch}}

demonstrate the abbreviated syntax and are also only matched in state \texttt{YYINITIAL}.
\label{ExampleRulesYYbegin}

Of these rules, one may be of special interest:

\verb+\"  { + \texttt{\htmlref{string.setLength(0); yybegin(STRING);}{CodeRulesBunch}}\verb+ }+

If the scanner matches a double quote in state \texttt{YYINITIAL} we
have recognised the start of a string literal. Therefore we clear our \texttt{StringBuffer}
that will hold the content of this string literal and tell the scanner
with \texttt{yybegin(STRING)} to switch into the lexical state \texttt{STRING}.
Because we do not yet return a value to the parser, our scanner proceeds
immediately.
\label{ExampleRulesYYtext}

In lexical state \texttt{STRING} another
rule demonstrates how to refer to the input that has been matched:

\verb-[^\n\r\"]+  { - \texttt{\htmlref{string.append( yytext() );}{CodeRulesYYtext}}\verb+ }+

The expression \verb-[^\n\r\"]+- matches
all characters in the input up to the next backslash (indicating an
escape sequence such as \verb+\n+), double quote (indicating the end
of the string), or line terminator (which must not occur in a string literal).
The matched region of the input is referred to with \texttt{\htmlref{yytext()}{CodeRulesYYtext}}
and appended to the content of the string literal parsed so far.
\label{ExampleRuleLast}

The last lexical rule in the example specification
is used as an error fallback. It matches any character in any state that
has not been matched by another rule. It doesn't conflict with any other
rule because it has the least priority (because it's the last rule) and
because it matches only one character (so it can't have longest match
precedence over any other rule).
 
\subsection{How to get it going}
\begin{itemize}
\item
Install JFlex (see section \ref{Installing} \htmlref{\emph{Installing JFlex}}{Installing})

\item
If you have written your specification file (or chosen one from the \texttt{examples}
directory), save it (say under the name \texttt{java-lang.flex}).

\item
Run JFlex with

\texttt{jflex java-lang.flex}

\item
JFlex should then report some progress messages about generating the scanner
and write the generated code to the directory of your specification file.

\item
Compile the generated \texttt{.java} file and your own classes. (If you
use CUP, generate your parser classes first)

\item
That's it.
\end{itemize}


\section{Lexical Specifications\label{Specifications}}
As shown above, a lexical specification file for JFlex consists of three
parts divided by a single line starting with \texttt{\%\%}:

\texttt{\htmlref{UserCode}{SpecUsercode}}\\
\texttt{\%\%}\\
\texttt{\htmlref{Options and declarations}{SpecOptions}}\\
\texttt{\%\%}\\
\texttt{\htmlref{Lexical rules}{LexRules}}

In all parts of the specification comments of the form
\texttt{/* comment text */} and the Java style end of line comments starting with \texttt{//}
are permitted. JFlex comments do nest - so the number of \texttt{/*} and \texttt{*/} 
should be balanced.

\subsection{User code\label{SpecUsercode}}
The first part contains user code that is copied verbatim into the beginning
of the source file of the generated lexer before the scanner class is declared. 
As shown in the example above, this is the place to put \texttt{package} 
declarations and \texttt{import}
statements. It is possible, but not considered good Java programming
style to put own helper classes (such as token classes) in this section.
They should get their own \texttt{.java} file instead.

\subsection{Options and declarations\label{SpecOptions}}
The second part of the lexical specification contains \htmlref{options}{SpecOptDirectives}
to customise your generated lexer (JFlex directives and Java code to include in
different parts of the lexer), declarations of \htmlref{lexical states}{StateDecl} and 
\htmlref{macro definitions}{MacroDefs} for use in the third section 
\htmlref{``Lexical rules''}{LexRules} of the lexical specification file.
\label{SpecOptDirectives}

Each JFlex directive must be situated at the beginning of  a line
and starts with the \texttt{\%} character. Directives that have one or
more parameters are described as follows:

\texttt{\%class "classname"}

means that you start a line with \texttt{\%class} followed by a space followed
by the name of the class for the generated scanner (the double quotes are
\emph{not} to be entered, see the \htmlref{example specification}{CodeOptions} in
section \ref{CodeOptions}).
 
\subsubsection{Class options and user class code\label{ClassOptions}}
These options regard name, constructor, API, and related parts of the
generated scanner class.
\begin{itemize}
\item
{\bf \texttt{\%class "classname"}}

Tells JFlex to give the generated class the name "\texttt{classname}" and to
write the generated code to a file "\texttt{classname.java}". If the 
\texttt{-d <directory>} command line option is not used, the code
will be written to the directory where the specification file resides. If
no \texttt{\%class} directive is present in the specification, the generated
class will get the name "\texttt{Yylex}" and will be written to a file
"\texttt{Yylex.java}". There should be only one \texttt{\%class} directive
in a specification.
 
\item
{\bf \texttt{\%implements "interface 1"[, "interface 2", ..]}}

Makes the generated class implement the specified interfaces. If more than
one \texttt{\%imple\-ments} directive is present, all the specified interfaces 
will be implemented.
 
\item
{\bf \texttt{\%extends "classname"}}

Makes the generated class a subclass of the class ``\texttt{classname}''.
There should be only one \texttt{\%extends} directive in a specification.

\item
{\bf \texttt{\%public}}

Makes the generated class public (the class is only accessible in its
own package by default).

\item
{\bf \texttt{\%final}}

Makes the generated class final.

\item
{\bf \texttt{\%abstract}}

Makes the generated class abstract.
 
\item
{\bf \texttt{\%apiprivate}}

Makes all generated methods and fields of the class
private. Exceptions are the constructor, user code in the
specification, and, if \verb+%cup+ is present, the method
\texttt{next\_token}. All occurrences of
\texttt{" public "} (one space character before and after \texttt{public})
in the skeleton file are replaced by
\texttt{" private "} (even if a user-specified skeleton is used).
Access to the generated class is expected to be mediated by user class
code (see next switch).


\item
{\bf \verb+%{+}\\ {\bf \texttt{...}}\\ {\bf \verb+%}+}

The code enclosed in \verb+%{+ and \verb+%}+ is copied verbatim
into the generated class. Here you can define your own member variables
and functions in the generated scanner. Like all options, both \verb+%{+
and \verb+%}+ must start a line in the specification. If more than one
class code directive \verb+%{...%}+ is present, the code is concatenated
in order of appearance in the specification.
 
\item
{\bf \verb+%init{+}\\
{\bf \texttt{...}}\\
{\bf \verb+%init}+}

The code enclosed in \verb+%init{+ and \verb+%init}+ is copied
verbatim into the constructor of the generated class. Here, member
variables declared in the \verb+%{...%}+ directive can be initialised.
If more than one initialiser option is present, the code is concatenated
in order of appearance in the specification.
 
\item
{\bf \verb+%initthrow{+}\\
{\bf \texttt{"exception1"[, "exception2", ...]}}\\
{\bf \verb+%initthrow}+}

or (on a single line) just

{\bf \texttt{\%initthrow "exception1" [,  "exception2", ...]}}

Causes the specified exceptions to be declared in the \texttt{throws}
clause of the constructor. If more than one \verb+%initthrow{+ \texttt{...} \verb+%initthrow}+
directive is present in the specification, all specified exceptions will
be declared.

\item
{\bf \texttt{\%ctorarg "type" "ident"}}

Adds the specified argument to the constructors of the generated scanner. 
If more than one such directive is present, the arguments are added in order
of occurrence in the specification. Note that this option conflicts with
the \verb+%standalone+ and \verb+%debug+ directives, because there is no
sensible default that can be created automatically for such parameters 
in the generated \texttt{main} methods. JFlex will warn in this case and
generate an additional default constructor without these parameters and without user init code (which might potentially refer to the parameters).

\item 
{\bf \texttt{\%scanerror "exception"}}

Causes the generated scanner to throw an instance of the specified
exception in case of an internal error (default is
\texttt{java.lang.Error}).  Note that this exception is only for
internal scanner errors. With usual specifications it should never
occur (i.e.~if there is an error fallback rule in the specification
and only the documented scanner API is used).

\item
{\bf \texttt{\%buffer "size"}}

Set the initial size of the scan buffer to the specified value 
(decimal, in bytes). The default value is 16384.

\item 
{\bf \texttt{\%include "filename"}}

Replaces the \texttt{\%include} verbatim by the specified file. This
feature is still experimental. It works, but error reporting can be
strange if a syntax error occurs on the last token in the included
file.

\end{itemize}

\subsubsection{Scanning method\label{ScanningMethod}}
This section shows how the scanning method can be customised. You can redefine
the name and return type of the method and it is possible to declare
exceptions that may be thrown in one of the actions of the specification.
If no return type is specified, the scanning method will be declared as
returning values of class \texttt{Yytoken}.
\begin{itemize}
\item
{\bf \texttt{\%function "name"}}

Causes the scanning method to get the specified name. If no \texttt{\%function}
directive is present in the specification, the scanning method gets the
name ``\texttt{yylex}''. This directive overrides settings of the 
\texttt{\htmlref{\%cup}{CupMode}} switch. Please note that the default name
of the scanning method with the \texttt{\htmlref{\%cup}{CupMode}} switch is
\texttt{next\_token}. Overriding this name might lead to the generated scanner
being implicitly declared as \texttt{abstract}, because it does not provide
the method \texttt{next\_token} of the interface \texttt{java\_cup.runtime.Scanner}.
It is of course possible to provide a dummy implementation of that method
in the class code section if you still want to override the function name.
 
\item
{\bf \texttt{\%integer}}\\
{\bf \texttt{\%int}}

Both cause the scanning method to be declared as of Java type \texttt{int}.
Actions in the specification can then return \texttt{int} values as tokens.
The default end of file value under this setting is \texttt{YYEOF}, which is a \texttt{public
static final int} member of the generated class.
 
\item
{\bf \texttt{\%intwrap}}

Causes the scanning method to be declared as of the Java wrapper type
\texttt{Integer}. Actions in the specification can then return \texttt{Integer}
values as tokens. The default end of file value under this setting is \texttt{null}.
 
\item
{\bf \texttt{\%type "typename"}}

Causes the scanning method to be declared as returning values of the specified type.
Actions in the specification can then return values of \texttt{typename}
as tokens. The default end of file value under this setting is \texttt{null}.
If \texttt{typename} is not a subclass of \texttt{java.lang.Object},
you should specify another end of file value using the
\htmlref{\texttt{\%eofval\bl} \texttt{...} \texttt{\%eofval\br}}{eofval}
directive or the \htmlref{\texttt{<<EOF>>} rule}{EOFRule}. 
The \texttt{\%type} directive overrides settings of the 
\texttt{\htmlref{\%cup}{CupMode}} switch.
 
\item
{\bf \verb+%yylexthrow{+}\\
{\bf \texttt{"exception1"[, "exception2", ... ]}}\\
{\bf \verb+%yylexthrow}+}

or (on a single line) just

{\bf \texttt{\%yylexthrow "exception1" [,  "exception2", ...]}}

The exceptions listed inside \verb+%yylexthrow{+ \texttt{...} \verb+%yylexthrow}+ 
will be declared in the throws clause of the scanning method. If there is
more than one \verb+%yylexthrow{+ \texttt{...} \verb+%yylexthrow}+ clause in
the specification, all specified exceptions will be declared.
\end{itemize}

\subsubsection{The end of file\label{EOF}}
There is always a default value that the scanning method will return when
the end of file has been reached. You may however define a specific value
to return and a specific piece of code that should be executed when the
end of file is reached.

The default end of file value depends on the return type of the scanning method:
\begin{itemize}
\item
For {\bf \texttt{\%integer}}, the scanning method will return the value
{\bf \texttt{YYEOF}}, which is a \texttt{public static final int} member
of the generated class.

\item
For {\bf \texttt{\%intwrap}}, 
\item
no specified type at all, or a 
\item
user defined type, declared using {\bf \texttt{\%type}}, the value is {\bf \texttt{null}}.

\item
In CUP compatibility mode, using {\bf \texttt{\%cup}}, the value is 

{\bf \texttt{new java\_cup.runtime.Symbol(sym.EOF)}}
\end{itemize}

User values and code to be executed at the end of file can be defined using these directives:
\begin{itemize}
\label{eofval}
\item
{\bf \verb+%eofval{+}\\
{\bf \texttt{...}}\\
{\bf \verb+%eofval}+}

The code included in \verb+%eofval{+ \texttt{...} \verb+%eofval}+ will
be copied verbatim into the scanning method and will be executed {\trit each time} 
when the end of file is reached (this is possible when
the scanning method is called again after the end of file has been
reached). The code should return the value that indicates the end of
file to the parser.  There should be only one \verb+%eofval{+ 
\texttt{...} \verb+%eofval}+ clause in the specification.  
The \verb+%eofval{ ... %eofval}+ directive overrides settings of the 
\texttt{\htmlref{\%cup}{CupMode}} switch and \texttt{\htmlref{\%byaccj}{YaccMode}} switch. 
As of version 1.2 JFlex provides
a more readable way to specify the end of file value using the
\htmlref{\texttt{<<EOF>>} rule}{EOFRule} (see also section \ref{EOFRule}).

\item\label{eof}
  {\bf \verb+%eof{+}\\
  {\bf \texttt{...}}\\
  {\bf \verb+%eof}+} 
  
  The code included in \verb+%{eof ... %eof}+ will be executed
  exactly once, when the end of file is reached. The code is included
  inside a method \texttt{void yy\_do\_eof()} and should not return any
  value (use \verb+%eofval{...%eofval}+ or 
  \htmlref{\texttt{<<EOF>>}}{EOFRule} for this purpose). If more than one 
  end of file code directive is present, the code will be concatenated
  in order of appearance in the specification.
 

\item
  {\bf \verb+%eofthrow{+}\\
  {\bf \texttt{"exception1"[,"exception2", ... ]}}\\
  {\bf \verb+%eofthrow}+}

  or (on a single line) just

  {\bf \texttt{\%eofthrow "exception1" [,  "exception2", ...]}}
  
  The exceptions listed inside \verb+%eofthrow{...%eofthrow}+ will
  be declared in the throws clause of the method \texttt{yy\_do\_eof()}
  (see \htmlref{\texttt{\%eof}}{eof} for more on that method).
  If there is more than one \verb+%eofthrow{...%eofthrow}+ clause
  in the specification, all specified exceptions will be declared.
 

\label{eofclose}
\item{\bf \texttt{\%eofclose}}
  
  Causes JFlex to close the input stream at the end of file. The code
  \texttt{yyclose()} is appended to the method \texttt{yy\_do\_eof()}
  (together with the code specified in \verb+%eof{...%eof}+) and
  the exception \texttt{java.io.IOException} is declared in the throws
  clause of this method (together with those of 
  \verb+%eofthrow{...%eofthrow}+)


\item{\bf \texttt{\%eofclose false}}
  
  Turns the effect of \texttt{\%eofclose} off again (e.g. in case closing of
  input stream is not wanted after \texttt{\%cup}).

\end{itemize}

\subsubsection{Standalone scanners\label{Standalone}}
\begin{itemize}
\item
{\bf \texttt{\%debug}}

Creates a main function in the generated class that expects the name
of an input file on the command line and then runs the scanner on this
input file by printing information about each returned token to the Java 
console until the end of file is reached. The information includes:
line number (if line counting is enabled), column (if column counting is enabled),
the matched text, and the executed action (with line number in the specification).
 
\item
{\bf \texttt{\%standalone}}

Creates a main function in the generated class that expects the name
of an input file on the command line and then runs the scanner on this
input file. The values returned by the scanner are ignored, but any unmatched
text is printed to the Java console instead (as the C/C++ tool flex does, if
run as standalone program). To avoid having to use an extra token class, the
scanning method will be declared as having default type \texttt{int}, not \texttt{YYtoken}
(if there isn't any other type explicitly specified).
This is in most cases irrelevant, but could be useful to know when making
another scanner standalone for some purpose. You should also consider using
the \texttt{\%debug} directive, if you just want to be able to run the scanner
without a parser attached for testing etc.

\end{itemize}


\subsubsection{CUP compatibility\label{CupMode}}
You may also want to read section \ref{CUPWork} \htmlref{\emph{JFlex and CUP}}{CUPWork}
if you are interested in how to interface your generated
scanner with CUP.
\begin{itemize}
\item
{\bf \texttt{\%cup}}

The \texttt{\%cup} directive enables the CUP compatibility mode and is equivalent
to the following set of directives:

\begin{verbatim}
%implements java_cup.runtime.Scanner
%function next_token
%type java_cup.runtime.Symbol
%eofval{
  return new java_cup.runtime.Symbol(<CUPSYM>.EOF);
%eofval}
%eofclose
\end{verbatim}

The value of \texttt{<CUPSYM>} defaults to \texttt{sym} and can be
changed with the \texttt{\%cupsym} directive. In JLex compatibility
mode (\texttt{--jlex} switch on the command line), \texttt{\%eofclose}
will not be turned on.

\item
{\bf \texttt{\%cupsym "classname"}}

Customises the name of the CUP generated class/interface 
containing the names of terminal tokens. Default is \texttt{sym}.
The directive should not be used after \texttt{\%cup}, but before.

\item
{\bf \texttt{\%cupdebug}}

Creates a main function in the generated class that expects the name
of an input file on the command line and then runs the scanner on this
input file. Prints line, column, matched text, and CUP symbol name for
each returned token to standard out.

\end{itemize}

\subsubsection{BYacc/J compatibility\label{YaccMode}}
You may also want to read section \ref{YaccWork} \htmlref{\emph{JFlex and BYacc/J}}{YaccWork}
if you are interested in how to interface your generated
scanner with Byacc/J.
\begin{itemize}
\item
{\bf \texttt{\%byacc}}

The \texttt{\%byacc} directive enables the BYacc/J compatibility mode and is equivalent
to the following set of directives:

\begin{verbatim}
%integer
%eofval{
  return 0;
%eofval}
%eofclose
\end{verbatim}

\end{itemize}

\subsubsection{Code generation\label{CodeGeneration}}
The following options define what kind of lexical analyser code JFlex
will produce. \texttt{\%pack} is the default setting and will be used,
when no code generation method is specified.
 
\begin{itemize}
\item 
  {\bf \tt \%switch}

  {\bf The \texttt{\%switch} code generation method is deprecated and 
  will be removed in JFlex \nextver.} With \texttt{\%switch} JFlex will
  generate a scanner that has the DFA hard coded into a nested switch
  statement. This method gives a good deal of compression in terms of
  the size of the compiled \texttt{.class} file while still providing
  very good performance. If your scanner gets to big though (say more
  than about 200 states) performance may vastly degenerate and you
  should consider using one of the \texttt{\%table} or \texttt{\%pack}
  directives. If your scanner gets even bigger (about 300 states),
  the Java compiler \texttt{javac} could produce corrupted code, that
  will crash when executed or will give you an \texttt{java.lang.VerifyError}
  when checked by the virtual machine. This is due to the size
  limitation of 64 KB of Java methods as described in the Java Virtual
  Machine Specification \cite{MachineSpec}. In this case you will be
  forced to use the \texttt{\%pack} directive, since \texttt{\%switch}
  usually provides more compression of the DFA table than the
  \texttt{\%table} directive.

 
\item
  {\bf \texttt{\%table}}
  
  {\bf The \texttt{\%table} code generation method is deprecated and 
  will be removed in JFlex \nextver.} The \texttt{\%table} direction causes
  JFlex to produce a classical table driven scanner that encodes its
  DFA table in an array.  In this mode, JFlex only does a small amount
  of table compression (see \cite{ParseTable}, \cite{SparseTable}, 
  \cite{Aho} and \cite{Maurer} for more details on the matter of table
  compression) and uses the same method that JLex did up to version
  1.2.1. See section \ref{performance} \htmlref{performance}{performance}
  of this manual to compare these methods. The same reason as above
  (64 KB size limitation of methods) causes the same problem, when the
  scanner gets too big. This is, because the virtual machine treats
  static initialisers of arrays as normal methods. You will in this
  case again be forced to use the \texttt{\%pack} directive to avoid
  the problem.
 
\item
  {\bf \texttt{\%pack}}
  
  \texttt{\%pack} causes JFlex to compress the generated DFA table and to
  store it in one or more string literals. JFlex takes care that the
  strings are not longer than permitted by the class file format.
  The strings have to be unpacked when
  the first scanner object is created and initialised.
  After unpacking, the DFA table is exactly the
  same as with option \texttt{\%table} --- the only extra work to be done
  at runtime is the unpacking process which is fast (not noticeable
  in normal cases). It is in time complexity proportional to the
  size of the expanded DFA table, and it is static,
  i.e.\ it is done only once per scanner class --- no matter
  how often it is instantiated.  Again, see section
  \ref{performance} \htmlref{performance}{performance}
  on the performance of these scanners
  With \texttt{\%pack}, there should be practically no
  limitation to the size of the scanner. \texttt{\%pack} is the default
  setting and will be used when no code generation method is specified.
\end{itemize}

\subsubsection{Character sets\label{CharacterSets}}
\begin{itemize}
\item
{\bf \texttt{\%7bit}}

Causes the generated scanner to use an 7 bit input character set (character
codes 0-127). If an input character with a code greater than 127 is
encountered in an input at runtime, the scanner will throw an \texttt{ArrayIndexOutofBoundsException}.
Not only because of this, you should consider using the \texttt{\%unicode} directive. 
See also section \ref{sec:encodings} for information about character encodings. This is the default in JLex compatibility mode.
 
\item
{\bf \texttt{\%full}}\\
{\bf \texttt{\%8bit}}

Both options cause the generated scanner to use an 8 bit input character
set (character codes 0-255). If an input character with a code greater
than 255 is encountered in an input at runtime, the scanner will throw
an \texttt{ArrayIndexOutofBoundsException}. Note that even if your platform
uses only one byte per character, the Unicode value of a character may
still be greater than 255. If you are scanning text files, you should
consider using the \texttt{\%unicode} directive. See also section \ref{sec:encodings}
for more information about character encodings.
 
\item
{\bf \texttt{\%unicode}}\\
{\bf \texttt{\%16bit}}

Both options cause the generated scanner to use the BMP (Basic Multilingual
Plane) of the Unicode input character set that Java supports natively
(character code points 0-65535). JFlex \ver does not yet support supplementary
characters above the BMP; support is planned for JFlex \nextver. There will be
no runtime overflow when using this set of input characters.
\texttt{\%unicode} does not mean that the scanner will read two bytes at a
time. What is read and what constitutes a character depends on the runtime
platform. See also section \ref{sec:encodings} for more information about
character encodings. This is the default unless the JLex compatibility mode
is used (command line option \texttt{--jlex}).

\label{caseless} \item {\bf \texttt{\%caseless}}\\ {\bf \texttt{\%ignorecase}}

This option causes JFlex to handle all characters and strings in the
specification as if they were specified in both uppercase and lowercase form.
This enables an easy way to specify a scanner for a language with case
insensitive keywords. The string "\texttt{break}" in a specification is for
instance handled like the expression \texttt{([bB][rR][eE][aA][kK])}. The
\texttt{\%caseless} option does not change the matched text and does not
affect character classes. So \texttt{[a]} still only matches the character
\texttt{a} and not \texttt{A}, too. Which letters are uppercase and which
lowercase letters, is defined by the Unicode standard. In JLex compatibility
mode (\texttt{--jlex} switch on the command line), \texttt{\%caseless}
and \texttt{\%ignorecase} also affect character classes.

 \end{itemize}
\subsubsection{Line, character and column counting\label{Counting}}
\begin{itemize}
\item
{\bf \texttt{\%char}}

Turns character counting on. The \texttt{int} member variable \texttt{yychar}
contains the number of characters (starting with 0) from the beginning
of input to the beginning of the current token.
 
\item
{\bf \texttt{\%line}}

Turns line counting on. The \texttt{int} member variable \texttt{yyline}
contains the number of lines (starting with 0) from the beginning of input
to the beginning of the current token.
 
\item
{\bf \texttt{\%column}}

Turns column counting on. The \texttt{int} member variable \texttt{yycolumn}
contains the number of characters (starting with 0) from the beginning
of the current line to the beginning of the current token.

 \end{itemize}

\subsubsection{Obsolete JLex options\label{Obsolete}}
\begin{itemize}
\item
{\bf \texttt{\%notunix}}

This JLex option is obsolete in JFlex but still recognised as valid directive.
It used to switch between Windows and Unix kind of line terminators (\verb+\r\n+
and \verb+\n+) for the \texttt{\$} operator in regular expressions. JFlex
always recognises both styles of platform dependent line terminators.
 
\item
{\bf \texttt{\%yyeof}}

This JLex option is obsolete in JFlex but still recognised as valid directive.
In JLex it declares a public member constant \texttt{YYEOF}. JFlex declares it in any case.
\end{itemize}

\subsubsection{State declarations\label{StateDecl}}
State declarations have the following form:

\texttt{\%s[tate] "state identifier" [, "state identifier", ... ]} for inclusive or\\
\texttt{\%x[state] "state identifier" [, "state identifier", ... ]} for exclusive states

There may be more than one line of state declarations, each starting with
\texttt{\%state} or \texttt{\%xstate} (the first character is sufficient,
\texttt{\%s} and \texttt{\%x} works, too). State identifiers are letters followed 
by a sequence of letters, digits or underscores. State identifiers can be separated 
by white-space or comma.

The sequence

\texttt{\%state STATE1}\\
\texttt{\%xstate STATE3, XYZ, STATE\_10}\\
\texttt{\%state ABC STATE5}

declares the set of identifiers \texttt{{STATE1, STATE3, XYZ,
    STATE\_10, ABC, STATE5}} as lexical states, \texttt{STATE1}, \texttt{ABC}, \texttt{STATE5}
as inclusive, and \texttt{STATE3}, \texttt{XYZ}, \texttt{STATE\_10} as exclusive. 
See also section
\ref{HowMatched} on the way lexical states influence how the input is
matched.
 
\subsubsection{Macro definitions\label{MacroDefs}}
A macro definition has the form

\texttt{macroidentifier = regular expression}

That means, a macro definition is a macro identifier (letter followed
by a sequence of letters, digits or underscores), that can later be
used to reference the macro, followed by optional white-space, followed
by an "\texttt{=}", followed by optional white-space, followed by a
regular expression (see section \ref{LexRules} \htmlref{\emph{lexical
    rules}}{LexRules} for more information about regular expressions).

The regular expression on the right hand side must be well formed and
must not contain the \verb+^+, \texttt{/} or \texttt{\$} operators. {\bf Differently
to JLex, macros are not just pieces of text that are expanded by copying}
- they are parsed and must be well formed.

{\bf This is a feature.} It eliminates some very hard to find bugs in
lexical specifications (such like not having parentheses around more
complicated macros - which is not necessary with JFlex).  See section
\ref{Porting} \htmlref{\emph{Porting from JLex}}{Porting} for more
details on the problems of JLex style macros.

Since it is allowed to have macro usages in macro definitions, it is
possible to use a grammar like notation to specify the desired lexical
structure. Macros however remain just abbreviations of the regular expressions
they represent. They are not non terminals of a grammar and cannot be used
recursively in any way. JFlex detects cycles in macro definitions and reports
them at generation time. JFlex also warns you about macros that have been
defined but never used in the ``lexical rules'' section of the specification.
  
\subsection{Lexical rules\label{LexRules}}
The ``lexical rules'' section of a JFlex specification contains a set of
regular expressions and actions (Java code) that are executed when the
scanner matches the associated regular expression.

The \texttt{\%include} directive may be used in this section to include
lexical rules from a separate file.  The directive will be replaced verbatim
by the contents of the specified file.
 
\subsubsection{Syntax\label{Grammar}}
The syntax of the "lexical rules" section is described by the following
EBNF grammar (terminal symbols are enclosed in 'quotes'):

\begin{verbatim}
LexicalRules ::= (Include|Rule)+
Include      ::= '%include' (' '|'\t'|'\b')+ File
Rule         ::= [StateList] ['^'] RegExp [LookAhead] Action 
               | [StateList] '<<EOF>>' Action
               | StateGroup 
StateGroup   ::= StateList '{' Rule+ '}' 
StateList    ::= '<' Identifier (',' Identifier)* '>' 
LookAhead    ::= '$' | '/' RegExp
Action       ::= '{' JavaCode '}' | '|'

RegExp       ::= RegExp '|' RegExp 
               | RegExp RegExp 
               | '(' RegExp ')'
               | ('!'|'~') RegExp
               | RegExp ('*'|'+'|'?')
               | RegExp "{" Number ["," Number] "}" 
               | CharClass
               | PredefinedClass 
               | MacroUsage 
               | '"' StringCharacter+ '"' 
               | Character 

CharClass    ::= '[' ['^'] CharClassContent* ']'
               | '[' ['^'] CharClassContent+ 
                     CharClassOperator CharClassContent+ ']'
                 
CharClassContent    ::= CharClass | Character |
                        Character'-'Character | 
                        MacroUsage | PredefinedClass

CharClassOperator   ::= '||' | '&&' | '--' | '~~'

MacroUsage          ::= '{' Identifier '}'

PredefinedClass     ::= '[:jletter:]' 
                      | '[:jletterdigit:]' 
                      | '[:letter:]' 
                      | '[:digit:]'
                      | '[:uppercase:]' 
                      | '[:lowercase:]'
                      | '\d' | '\D'
                      | '\s' | '\S'
                      | '\w' | '\W'
                      | '\p{' UnicodePropertySpec '}'
                      | '\P{' UnicodePropertySpec '}'
                      | '\R'
                      | '.'          
                            
UnicodePropertySpec ::= BinaryProperty | 
                        EnumeratedProperty (':' | '=') PropertyValue

BinaryProperty      ::= Identifier

EnumeratedProperty  ::= Identifier

PropertyValue       ::= Identifier
\end{verbatim}
%$

\label{Terminals}
The grammar uses the following terminal symbols:   
\begin{itemize}
\item \texttt{File}\\
  a file name, either absolute or relative to the directory containing 
  the lexical specification.

\item \texttt{JavaCode}\\
  a sequence of {\trit \texttt{BlockStatements}} as described in the Java
  Language Specification \cite{LangSpec}, section 14.2.

\item \texttt{Number}\\
  a non negative decimal integer.
  
\item \texttt{Identifier}\\
  a letter \verb+[a-zA-Z]+ followed by a sequence of zero or more
  letters, digits or underscores \verb+[a-zA-Z0-9_]+

\item \texttt{Character}\\
  an escape sequence or any unicode character that is not one of these
  meta characters:
  \verb:  |  (  )  {  }  [  ]  < >  \  .  *  +  ?  ^  $  / . " ~ !:
%$

\item \texttt{StringCharacter}\\   
  an escape sequence or any unicode character that is not one of these
  meta characters:
  \verb:  \  ":

\item
  An escape sequence

  \begin{itemize}
  \item
    \verb+\n+  \verb+\r+  \verb+\t+  \verb+\f+  \verb+\b+
  \item
    a \verb+\x+ followed by two hexadecimal digits \texttt{[a-fA-F0-9]} (denoting
    a standard ASCII escape sequence),
    
  \item
    a \verb+\u+ followed by four hexadecimal digits \texttt{[a-fA-F0-9]}
    (denoting a unicode escape sequence),

  \item
    a backslash followed by a three digit octal number from 000 to 377 (denoting
    a standard ASCII escape sequence), or

  \item
    a backslash followed by any other unicode character that stands for this
    character.

  \end{itemize}
  
\end{itemize}

Please note that the \verb+\n+ escape sequence stands for the ASCII
LF character - not for the end of line. If you would like to match the
line terminator, you should use the expression \verb+\r|\n|\r\n+ if you want
the Java conventions, or \verb+\r\n|[\r\n\u2028\u2029\u000B\u000C\u0085]+
(provided as predefined class \verb+\R+) if you want to be fully Unicode
compliant (see also \cite{unicode_rep}).

As of version 1.1 of JFlex the white-space characters \texttt{" "}
(space) and \verb+"\t"+ (tab) can be used to improve the readability of
regular expressions. They will be ignored by JFlex. In character
classes and strings however, white-space characters keep standing for
themselves (so the string \texttt{" "} still matches exactly one space
character and \verb+[ \n]+ still matches an ASCII LF or a space
character).

JFlex applies the following standard operator precedences in regular
expression (from highest to lowest):

\begin{itemize}
\item
unary postfix operators (\verb-'*', '+', '?', {n}, {n,m}-)

\item
unary prefix operators (\verb-'!', '~'-)

\item
concatenation (\texttt{RegExp::= RegExp Regexp})

\item
union (\verb-RegExp::= RegExp '|' RegExp-)
\end{itemize}

So the expression \verb+a | abc | !cd*+ for instance is parsed as 
\verb+(a|(abc)) | ((!c)(d*))+.

\subsubsection{Semantics\label{Semantics}}
This section gives an informal description of which text is matched by
a regular expression (i.e. an expression described by the \texttt{RegExp}
production of the grammar presented \htmlref{above}{Grammar}).

A regular expression that consists solely of
\begin{itemize}
\item
  a \texttt{Character} matches this character.

\item
  a character class \verb:'[...]': matches any character in that class.
  A \texttt{Character} is to be considered an element of a class, if it
  is listed in the class or if its code lies within a listed character
  range \texttt{Character'-'Character} or Macro or predefined character
  class. So \verb+[a0-3\n]+ for instance matches the characters

  \verb+a 0 1 2 3 \n+

  If the list of characters is empty (i.e.~just \verb+[]+), the expression
  matches nothing at all (the empty set), not even the empty string. This
  may be useful in combination with the negation operator \verb+'!'+.
  
  Character sets may be nested, e.g. \verb+[[[abc]d[e]]fg]+ is equivalent
  to \verb+[abcdefg]+.

  Supported character set operations:
    
  \begin{itemize}
  \item
    Union (\verb+||+), e.g. \verb+[[a-c]||[d-f]]+, equivalent to 
    \verb+[a-cd-f]+: this is the default character set operation when
    no operator is specified. 

  \item
    Intersection (\verb+&&+), e.g. \verb+[[a-f]&&[f-m]]+, equivalent to
    \verb+[f]+.
    
  \item
    Set difference (\verb+--+), e.g. \verb+[[a-z]--m]+, equivalent to
    \verb+[a-ln-z]+.
    
  \item
    Symmetric difference (\verb+~~+): the union of two classes minus their
    intersection. For instance \verb+[\p{Letter}~~\p{ASCII}]+ is equivalent
    to
    \verb+[[\p{Letter}||\p{ASCII}]--+
    \verb+[\p{Letter}&&\p{ASCII}]]+: the
    set of characters that are present in either \verb+\p{Letter}+ or
    in \verb+\p{ASCII}+, but not in both.

  \end{itemize}

\item 
  a negated character class \verb:'[^...]':  matches all characters not
  listed in the class. If the list of characters is empty (i.e. 
  \verb+[^]+), the expression matches any character of the input
  character set.
  
\item
  a string \texttt{'"' StringCharacter+ '"} \texttt{'} matches the exact
  text enclosed in double quotes. All meta characters but \verb+\+ and
  \texttt{"} lose their special meaning inside a string. See also the
  \htmlref{\texttt{\%ignorecase}}{caseless} switch.
    
\item
  a macro usage \verb+'{' Identifier '}'+ matches the input that is matched
  by the right hand side of the macro with name "\texttt{Identifier}".

\label{predefCharCl}
\item 
  a predefined character class matches any of
  the characters in that class. There are the following predefined character
  classes:

  \begin{itemize}
  \item
    These two predefined character classes are determined by Java functions of
    class \texttt{java}.\texttt{lang}.\texttt{Cha\-rac\-ter}:

    \begin{verbatim}
    [:jletter:]       isJavaIdentifierStart()
    [:jletterdigit:]  isJavaIdentifierPart()
    \end{verbatim}

  \item
    These four predefined character classes are equivalent to the following
    Unicode properties (described \htmlref{below}{unipropsyntax}):

    \begin{verbatim}
    [:letter:]     \p{Letter}
    [:digit:]      \p{Digit}
    [:uppercase:]  \p{Uppercase}
    [:lowercase:]  \p{Lowercase}
    \end{verbatim}

  \item
    These meta characters are equivalent to the following (sets of) Unicode
    Properties (described \htmlref{below}{unipropsyntax}):
    
    \begin{verbatim}
    \d  \p{Digit}
    \D  \P{Digit}
    \s  \p{Whitespace}
    \S  \P{Whitespace}
    \w  [\p{Alpha}\p{Digit}\p{Mark}
         \p{Connector Punctuation}\p{Join Control}]
    \W  [^\p{Alpha}\p{Digit}\p{Mark}
          \p{Connector Punctuation}\p{Join Control}]
    \end{verbatim}
    
  \label{unipropsyntax}
  \item
    Unicode Properties are character classes specified by each version of the
    Unicode Standard. JFlex supports a subset of all defined Properties for
    each supported Unicode version. To see the full list of supported
    Properties, give the \texttt{--uniprops <ver>} option on the JFlex
    command line, where \texttt{<ver>} is the Unicode version. Some
    Properties have aliases; JFlex recognizes all aliases for all supported
    properties. JFlex supports loose matching of Properties: case
    distinctions, whitespace, hyphens, and underscores are ignored.

    To refer to a Unicode Property, use the \verb+\p{...}+ syntax, e.g. the
    Greek Block can be referred to as \verb+\p{Block:Greek}+. To match the
    all characters not included in a property, use the \verb+\P{...}+ syntax
    (note that the '\verb+P+' is uppercase), e.g. to match all characters
    that are {\bf not} letters: \verb+\P{Letter}+.
     
    See UTS\#18 \cite{unicode_rep} for a description of and links to
    definitions of some supported Properties. UnicodeSet \cite{UnicodeSet}
    is an online utility to show the character sets corresponding to
    Unicode Properties and set operations on them, but only for the most
    recent Unicode version.

  \item
    Dot (\texttt{.}) matches \verb+[^\r\n\u2028\u2029\u000B\u000C\u0085]+.\\
    Use the \texttt{--legacydot} option to instead match \verb+[^\n]+.
    
  \item
    \verb+\R+ matches any newline: \verb+\r\n|[\r\n\u2028\u2029\u000B\u000C\u0085]+.
  
  \end{itemize}

\end{itemize}

    If \texttt{a} and \texttt{b} are regular expressions, then

    \begin{itemize}
     
          \item[\texttt{a | b}] (union) 

            is the regular expression that matches
            all input matched by \texttt{a} or by \texttt{b}.
          
          
          \item[\texttt{a b}] (concatenation) 
            
            is the regular expression
            that matches the input matched by \texttt{a} followed by the 
            input matched by \texttt{b}.
          
          
          \item[\texttt{a*}] (Kleene closure) 

            matches zero or more repetitions
            of the input matched by \texttt{a}
          
          
          \item[\texttt{a+}] (iteration)

          is equivalent to \texttt{aa*}
          
          
          \item[\texttt{a?}] (option)

          matches the empty input or the input matched
            by \texttt{a}

          \item[\texttt{!a}] (negation)

          matches everything but the strings matched by \texttt{a}. 
          Use with care: the construction of \verb+!a+ involves
          an additional, possibly exponential NFA to DFA transformation 
          on the NFA for \texttt{a}. Note that
          with negation and union you also have (by applying DeMorgan)
          intersection and set difference: the intersection of 
          \texttt{a} and \texttt{b} is \verb+!(!a|!b)+, the expression 
          that matches everything of \texttt{a} not matched by \texttt{b} is 
          \verb+!(!a|b)+

          \item[\texttt{\symbol{126}a}] (upto)

          matches everything up to (and including) the first occurrence of a text
          matched by \texttt{a}. The expression \verb-~a- is equivalent
          to \verb-!([^]* a [^]*) a-. A traditional C-style comment
          is matched by \verb-"/*" ~"*/"-
          
          \item[\texttt{a\bl n\br}] (repeat)          

            is equivalent to \texttt{n} times the concatenation of \texttt{a}.
            So \verb+a{4}+ for instance is equivalent to the expression \texttt{a a a a}.
            The decimal integer \texttt{n} must be positive.          
          
          \item[\texttt{a\bl n,m\br}]          
            is equivalent to at least \texttt{n} times and at most \texttt{m} times the 
            concatenation of \texttt{a}. So \verb+a{2,4}+ for instance is equivalent 
            to the expression \verb+a a a? a?+. Both \texttt{n} and \texttt{m} are non
            negative decimal integers and \texttt{m} must not be smaller than \texttt{n}.
          
          \item[\texttt{( a )}]
            matches the same input as \texttt{a}.
            
        
\end{itemize}
    
In a lexical rule, a regular expression \texttt{r} may be preceded by a
'\verb+^+' (the beginning of line operator). \texttt{r} is then
only matched at the beginning of a line in the input. A line begins
after each occurrence of \verb+\r|\n|\r\n|\u2028|\u2029|\u000B|\u000C|\u0085+ 
(see also \cite{unicode_rep}) and at the beginning of input.  
The preceding line terminator in the input is not consumed and can 
be matched by another rule.

In a lexical rule, a regular expression \texttt{r} may be followed by a
look-ahead expression. A look-ahead expression is either a '\texttt{\$}'
(the end of line operator) or a \verb+'/'+ followed by an arbitrary
regular expression. In both cases the look-ahead is not consumed and
not included in the matched text region, but it {\trit is} considered
while determining which rule has the longest match (see also 
\ref{HowMatched} \htmlref{\emph{How the input is matched}}{HowMatched}). 

In the '\texttt{\$}' case \texttt{r} is only matched at the end of a line in
the input. The end of a line is denoted by the regular expression
\verb+\r|\n|\r\n|\u2028|\u2029|\u000B|\u000C|\u0085+. 
So \verb+a$+ is equivalent to \verb+a / \r|\n|\r\n|\u2028|\u2029|\u000B|\u000C|\u0085+.
This is different to the situation described in \cite{unicode_rep}:
since in JFlex \verb+$+ is a true trailing context, the end of file
does {\bf not} count as end of line.

\label{trailingContext}
For arbitrary look-ahead (also called {\trit trailing context}) the 
expression is matched only when followed by input that matches the
trailing context. 

\label{EOFRule}
As of version 1.2, JFlex allows lex/flex style \texttt{<<EOF>>} rules in
lexical specifications. A rule
\begin{verbatim}
[StateList]  <<EOF>>    { some action code }
\end{verbatim}
is very similar to the \htmlref{\texttt{\%eofval} directive}{eofval} (section \ref{eofval}). 
The difference lies in the optional \texttt{StateList} that may precede the \texttt{<<EOF>>} rule. The
action code will only be executed when the end of file is read and the
scanner is currently in one of the lexical states listed in \texttt{StateList}. 
The same \texttt{StateGroup} (see section \ref{HowMatched} 
\htmlref{\emph{How the input is matched}}{HowMatched}) and precedence 
rules as in the ``normal'' rule case apply 
(i.e. if there is more than one \texttt{<<EOF>>} 
rule for a certain lexical state, the action of the one appearing 
earlier in the specification will be executed). \texttt{<<EOF>>} rules 
override settings of the \texttt{\%cup} and \texttt{\%byaccj} options and 
should not be mixed with the \texttt{\%eofval} directive.

An \texttt{Action} consists either of a piece of Java code enclosed in
curly braces or is the special \verb+|+ action. The \verb+|+ action is
an abbreviation for the action of the following expression.

Example:
\begin{verbatim}
expression1   |
expression2   |
expression3   { some action }
\end{verbatim}
is equivalent to the expanded form
\begin{verbatim}
expression1   { some action }
expression2   { some action }
expression3   { some action }
\end{verbatim}

They are useful when you work with trailing context expressions. The 
expression \texttt{a | (c / d) | b} is not syntactically legal, but can 
easily be expressed using the \verb+|+ action:
\begin{verbatim}
a       |
c / d   |
b       { some action }
\end{verbatim}

\subsubsection{How the input is matched\label{HowMatched}}
When consuming its input, the scanner determines the regular expression
that matches the longest portion of the input (longest match rule). If
there is more than one regular expression that matches the longest portion
of input (i.e. they all match the same input), the generated scanner chooses
the expression that appears first in the specification. After determining
the active regular expression, the associated action is executed. If there
is no matching regular expression, the scanner terminates the program with
an error message (if the \texttt{\%standalone} directive has been used, the
scanner prints the unmatched input to \texttt{java.lang.System.out} instead
and resumes scanning). 

Lexical states can be used to further restrict the set of regular expressions
that match the current input. 
 
\begin{itemize}
\item
A regular expression can only be matched when its associated set of lexical
states includes the currently active lexical state of the scanner or if
the set of associated lexical states is empty and the currently active lexical
state is inclusive. Exclusive and inclusive states only differ at this point: 
rules with an empty set of associated states.
 
\item
The currently active lexical state of the scanner can be changed from within
an action of a regular expression using the method \texttt{yybegin()}.
 

\item
The scanner starts in the inclusive lexical state 
\texttt{YYINITIAL}, which is always declared by default.
 

\item
The set of lexical states associated with a regular expression is 
the \texttt{StateList} that precedes the expression. If a rule is
contained in one or more \texttt{StateGroups}, then the states of
these are also associated with the rule, i.e.~they accumulate over
\texttt{StateGroups}.
  
Example:
\begin{verbatim}
%states A, B
%xstates C
%%
expr1                   { yybegin(A); action }
<YYINITIAL, A> expr2    { action }
<A> {
  expr3                 { action }
  <B,C> expr4           { action }
}
\end{verbatim}
The first line declares two (inclusive) lexical states \texttt{A} and \texttt{B},
the second line an exclusive lexical state \texttt{C}.
The default (inclusive) state \texttt{YYINITIAL} is always implicitly there and
doesn't need to be declared. The rule with \texttt{expr1} has no
states listed, and is thus matched in all states but the exclusive
ones, i.e.~\texttt{A}, \texttt{B}, and \texttt{YYINITIAL}. In its
action, the scanner is switched to state \texttt{A}. The second rule
\texttt{expr2} can only match when the scanner is in state
\texttt{YYINITIAL} or \texttt{A}. The rule \texttt{expr3} can only be
matched in state \texttt{A} and \texttt{expr4} in states \texttt{A}, \texttt{B},
and \texttt{C}.

\item
Lexical states are declared and used as Java \texttt{int} constants in
the generated class under the same name as they are used in the specification.
There is no guarantee that the values of these integer constants are
distinct. They are pointers into the generated DFA table, and if JFlex
recognises two states as lexically equivalent (if they are used with the
exact same set of regular expressions), then the two constants will get
the same value.
 
\end{itemize}

\subsubsection{The generated class}
JFlex generates exactly one file containing one class from the specification
(unless you have declared another class in the first specification section).

The generated class contains (among other things) the DFA tables, an input buffer, 
the lexical states of the specification, a constructor, and the scanning method
with the user supplied actions.

The name of the class is by default \texttt{Yylex}, it is customisable
with the \texttt{\%class} directive (see also section
\ref{ClassOptions}). The input buffer of the lexer is connected with an
input stream over the \texttt{java.io.Reader} object which is passed
to the lexer in the generated constructor. If you want to provide your
own constructor for the lexer, you should always call the generated
one in it to initialise the input buffer. The input buffer should not
be accessed directly, but only over the advertised API (see also
section \ref{ScannerMethods}). Its internal implementation may change
between releases or skeleton files without notice.

The main interface to the outside world is the generated scanning
method (default name \texttt{yylex}, default return type
\texttt{Yytoken}). Most of its aspects are customisable (name, return
type, declared exceptions etc., see also section
\ref{ScanningMethod}).  If it is called, it will consume input until
one of the expressions in the specification is matched or an error
occurs. If an expression is matched, the corresponding action is
executed. It may return a value of the specified return type (in which
case the scanning method returns with this value), or if it doesn't
return a value, the scanner resumes consuming input until the next
expression is matched. If the end of file is reached, the scanner
executes the EOF action, and (also upon each further call to the scanning
method) returns the specified EOF value (see also section \ref{EOF}).


\subsubsection{Scanner methods and fields accessible in actions (API)\label{ScannerMethods}}
Generated methods and member fields in JFlex scanners are prefixed
with \texttt{yy} to indicate that they are generated and to avoid name
conflicts with user code copied into the class. Since user code is
part of the same class, JFlex has no language means like the
\texttt{private} modifier to indicate which members and methods are
internal and which ones belong to the API. Instead, JFlex follows a
naming convention: everything starting with a \texttt{zz} prefix like
\texttt{zzStartRead} is to be considered internal and subject to
change without notice between JFlex releases. Methods and members of
the generated class that do not have a \texttt{zz} prefix like
\texttt{yycharat} belong to the API that the scanner class provides to
users in action code of the specification. They will remain stable
and supported between JFlex releases as long as possible.

Currently, the API consists of the following methods and member fields:
\begin{itemize}
\item \texttt{String yytext()}\\
  returns the matched input text region

\item \texttt{int yylength()}\\
  returns the length of the matched input text region (does not require
  a \texttt{String} object to be created)
  
\item \texttt{char yycharat(int pos)}\\
  returns the character at position \texttt{pos} from the matched text.
  It is equivalent to \texttt{yytext().charAt(pos)}, but faster.  {\tt
  pos} must be a value from \texttt{0} to \texttt{yylength()-1}.

\item \texttt{void yyclose()}\\
  closes the input stream. All subsequent calls to the scanning method will 
  return the end of file value    

\item \texttt{void yyreset(java.io.Reader reader)}\\
  closes the current input stream, and resets the scanner to read from
  a new input stream.  All internal variables are reset, the old input
  stream {\em cannot} be reused (content of the internal buffer is
  discarded and lost).  The lexical state is set to \texttt{YY\_INITIAL}.

\item \texttt{void yypushStream(java.io.Reader reader)}\\
 Stores the current input stream on a stack, and
 reads from a new stream. Lexical state, line,
 char, and column counting remain untouched.
 The current input stream can be restored with
 \texttt{yypopstream} (usually in an \texttt{<<EOF>>} action).
 
 A typical example for this are include files in
 style of the C pre-processor. The corresponding 
 JFlex specification could look somewhat like this:
\begin{verbatim}
"#include" {FILE}  { yypushStream(new FileReader(getFile(yytext()))); }
..
<<EOF>>        { if (yymoreStreams()) yypopStream(); else return EOF; }
\end{verbatim}

 This method is only available in the skeleton file
 \texttt{skeleton.nested}. You can find it in the 
 \texttt{src} directory of the JFlex distribution.

\item \texttt{void yypopStream()}\\
  Closes the current input stream and continues to
  read from the one on top of the stream stack.

  This method is only available in the skeleton file
  \texttt{skeleton.nested}. You can find it in the 
  \texttt{src} directory of the JFlex distribution.

\item \texttt{boolean yymoreStreams()}\\
  Returns true iff there are still streams for \texttt{yypopStream} 
  left to read from on the stream stack.

  This method is only available in the skeleton file
  \texttt{skeleton.nested}. You can find it in the 
  \texttt{src} directory of the JFlex distribution. 

\item \texttt{int yystate()}\\
  returns the current lexical state of the scanner.

\item \texttt{void yybegin(int lexicalState)}\\
  enters the lexical state \texttt{lexicalState}

\item \texttt{void yypushback(int number)}\\
  pushes \texttt{number} characters of the matched text back into the input stream. 
  They will be read again in the next call of the scanning method. 
  The number of characters to be read again must not be greater than the length
  of the matched text. The pushed back characters will not be included in \texttt{yylength()} and \texttt{yytext()}.
  Note that in Java strings are unchangeable, i.e.\ an action code like
  \begin{verbatim}
    String matched = yytext();
    yypushback(1);
    return matched;
  \end{verbatim}
  will return the whole matched text, while    
  \begin{verbatim}
    yypushback(1);
    return yytext();
  \end{verbatim}
  will return the matched text minus the last character.

\item\texttt{int yyline}\\
  contains the current line of input (starting with 0, only active with
  the \texttt{\htmlref{\%line}{Counting}} directive)

\item \texttt{int yychar}\\
  contains the current character count in the input (starting with 0,
  only active with the \texttt{\htmlref{\%char}{Counting}} directive)

\item \texttt{int yycolumn}\\
  contains the current column of the current line (starting with 0, only
  active with the \texttt{\htmlref{\%column}{Counting}} directive)

\end{itemize}

\section{Encodings, Platforms, and Unicode\label{sec:encodings}}

This section tries to shed some light on the issues of Unicode and
encodings, cross platform scanning, and how to deal with binary data.
My thanks go to Stephen Ostermiller for his input on this topic.

\subsection{The Problem\label{sec:howtoencoding}}

Before we dive straight into details, let's take a look at what the
problem is. The problem is Java's platform independence when you want
to use it. For scanners the interesting part about platform
independence is character encodings and how they are handled.

If a program reads a file from disk, it gets a stream of bytes.  In
earlier times, when the grass was green, and the world was much
simpler, everybody knew that the byte value 65 is, of course, an A.
It was no problem to see which bytes meant which characters (actually
these times never existed, but anyway).  The normal Latin alphabet
only has 26 characters, so 7 bits or 128 distinct values should surely
be enough to map them, even if you allow yourself the luxury of upper
and lower case.  Nowadays, things are different. The world suddenly
grew much larger, and all kinds of people wanted all kinds of special
characters, just because they use them in their language and writing.
This is were the mess starts. Since the 128 distinct values were
already filled up with other stuff, people began to use all 8 bits of
the byte, and extended the byte/character mappings to fit their need,
and of course everybody did it differently. Some people for instance
may have said ``let's use the value 213 for the German character {\"a}''.  Others
may have found that 213 should much rather mean {\'e}, because they didn't need
German and wrote French instead. As long as you use your program and
data files only on one platform, this is no problem, as all know what
means what, and everything gets used consistently.

Now Java comes into play, and wants to run everywhere (once written,
that is) and now there suddenly is a problem: how do I get the same
program to say {\"a} to a certain byte when it runs in Germany and maybe {\'e}
when it runs in France? And also the other way around: when I want to
say {\'e} on the screen, which byte value should I send to the operating
system?

Java's solution to this is to use Unicode internally.  Unicode aims to
be a superset of all known character sets and is therefore a perfect base
for encoding things that might get used all over the world. To make
things work correctly, you still have to know where you are and how to
map byte values to Unicode characters and vice versa, but the
important thing is, that this mapping is at least possible (you can
map Kanji characters to Unicode, but you cannot map them to ASCII or
iso-latin-1).

\subsection{Scanning text files\label{sec:howtotext}}

Scanning text files is the standard application for scanners like
JFlex. Therefore it should also be the most convenient one. Most times 
it is.

The following scenario works like a breeze:
You work on a platform X, write your lexer specification there, can
use any obscure Unicode character in it as you like, and compile the
program. Your users work on any platform Y (possibly but not 
necessarily something different from X), they write their input files
on Y and they run your program on Y. No problems.

Java does this as follows:
If you want to read anything in Java that is supposed to contain text, 
you use a \texttt{FileReader} or some \texttt{InputStream} together with
an \texttt{InputStreamReader}. \texttt{InputStreams} return the raw bytes, the 
\texttt{InputStreamReader} converts the bytes into Unicode characters with
the platform's default encoding. If a text file is produced on the
same platform, the platform's default encoding should do the mapping
correctly. Since JFlex also uses readers and Unicode internally, this
mechanism also works for the scanner specifications. If you write an
\texttt{A} in your text editor and the editor uses the platform's encoding (say \texttt{A} is 65), 
then Java translates this into the logical Unicode \texttt{A} internally. 
If a user writes an \texttt{A} on a completely different platform (say \texttt{A} is 237 there),
then Java also translates this into the logical Unicode \texttt{A} internally. Scanning
is performed after that translation and both match.

Note that because of this mapping from bytes to characters, you should always 
use the \texttt{\%unicode} switch in you lexer specification if you want to scan
text files. \texttt{\%8bit} may not be enough, even if
you know that your platform only uses one byte per character. The encoding
Cp1252 used on many Windows machines for instance knows 256 characters, but
the character {\'{}} with Cp1252 code \verb+\x92+ has the Unicode value \verb+\u2019+, which
is larger than 255 and which would make your scanner throw an 
\texttt{ArrayIndexOutOfBoundsException} if it is encountered.

So for the usual case you don't have to do anything but use the
\texttt{\%unicode} switch in your lexer specification.

Things may break when you produce a text file on platform X and
consume it on a different platform Y. Let's say you have a file
written on a Windows PC using the encoding Cp1252. Then you move
this file to a Linux PC with encoding ISO 8859-1 and there you want
to run your scanner on it. Java now thinks the file is encoded
in ISO 8859-1 (the platform's default encoding) while it really is 
encoded in Cp1252. For most characters
Cp1252 and ISO 8859-1 are the same, but for the byte values \verb+\x80+
to \verb+\x9f+ they disagree: ISO 8859-1 is undefined there. You can fix
the problem by telling Java explicitly which encoding to use. When
constructing the \texttt{InputStreamReader}, you can give the encoding
as argument. The line
\begin{center}
\texttt{Reader r = new InputStreamReader(input, "Cp1252"); }
\end{center}
will do the trick.

Of course the encoding to use can also come from the data itself:
for instance, when you scan an HTML page, it may have embedded 
information about its character encoding in the headers.

More information about encodings, which ones are supported, how
they are called, and how to set them may be found in the
official Java documentation in the chapter about 
internationalisation. 
The link 
\htmladdnormallink{\texttt{http://docs.oracle.com/javase/1.5.0/docs/guide/intl/}}{http://docs.oracle.com/javase/1.5.0/docs/guide/intl/}
leads to an online version of this for Oracle's JDK 1.5.

\subsection{Scanning binaries\label{sec:howtobinary}}

Scanning binaries is both easier and more difficult
than scanning text files. It's easier because you want
the raw bytes and not their meaning, i.e.~you don't want 
any translation.
It's more difficult because it's not so easy to get
``no translation'' when you use Java readers. 

The problem (for binaries) is that JFlex scanners are
designed to work on text. Therefore the interface is
the \texttt{Reader} class (there is a constructor
for \texttt{InputStream} instances, but it's just there
for convenience and wraps an \texttt{InputStreamReader}
around it to get characters, not bytes). 
You can still get a binary scanner when you write 
your own custom \texttt{InputStreamReader} class that
does explicitly no translation, but just copies
byte values to character codes instead. It sounds
quite easy, and actually it is no big deal, but there 
are a few little pitfalls on the way. In the scanner
specification you can only enter positive character
codes (for bytes that is \verb+\x00+
to \verb+\xFF+). Java's \texttt{byte} type on the other hand
is a signed 8 bit integer (-128 to 127), so you have to convert 
them properly in your custom \texttt{Reader}. Also, you should
take care when you write your lexer spec: if you
use text in there, it gets interpreted by an encoding
first, and what scanner you get as result might depend
on which platform you run JFlex on when you generate
the scanner (this is what you want for text, but for binaries it
gets in the way). If you are not sure, or if the development
platform might change, it's probably best to use character 
code escapes in all places, since they don't change their
meaning.

\section{Conformance with Unicode Regular Expressions UTS\#18\label{unicoderegexconformance}}

This section gives details about JFlex \ver's conformance with the
requirements for Basic Unicode Support Level 1 given in UTS\#18
\cite{unicode_rep}.

\subsubsection*{RL1.1 Hex Notation}

\begin{quote}
\emph{To meet this requirement, an implementation shall supply a mechanism
for specifying any Unicode code point (from U+0000 to U+10FFFF), using
the hexadecimal code point representation.}
\end{quote}

JFlex does not fully conform: although syntax is provided to express
values across the whole range, via \verb+\uXXXX+, where \verb+XXXX+ is
a 4-digit hex value, and \verb+\Uyyyyyy+, where \verb+yyyyyy+ is a 
6-digit hex value, JFlex only supports characters within the 16-bit
Basic Multilingual Plane, so the \verb+\Uyyyyyy+ syntax is not usable.

\subsubsection*{RL1.2 Properties}

\begin{quote}
\emph{To meet this requirement, an implementation shall provide at least a
minimal list of properties, consisting of the following: General\_Category,
Script and Script\_Extensions, Alphabetic, Uppercase, Lowercase, White\_Space,
Noncharacter\_Code\_Point, Default\_Ignorable\_Code\_Point, ANY, ASCII, ASSIGNED.}

\emph{The values for these properties must follow the Unicode definitions, 
and include the property and property value aliases from the UCD. Matching
of Binary, Enumerated, Catalog, and Name values, must follow the
Matching Rules from [UAX44].}
\end{quote}

JFlex conforms.  The minimal set of properties is supported, as well as
a few others. To see the full list of supported properties, use the JFlex
command line option \verb+--uniprops <ver>+, where \verb+<ver>+ is the
Unicode version.  Loose matching is performed: case distinctions, 
whitespace, underscores and hyphens in property names and values are
ignored.

\subsubsection*{RL1.2a Compatibility Properties}

\begin{quote}
\emph{To meet this requirement, an implementation shall provide the properties
listed in Annex C: Compatibility Properties, with the property values as
listed there. Such an implementation shall document whether it is using
the Standard Recommendation or POSIX-compatible properties.}
\end{quote}

JFlex does not fully conform. The Standard Recommendation version of the
Annex C Compatibility Properties are provided, with two exceptions:
\verb+\X+ Extended Grapheme Clusters; and \verb+\b+ Default Word
Boundaries.

\subsubsection*{RL1.3 Subtraction and Intersection}

\begin{quote}
\emph{To meet this requirement, an implementation shall supply mechanisms for
union, intersection and set-difference of Unicode sets.}
\end{quote}

JFlex conforms by providing these mechanisms, as well as symmetric
difference.

\subsubsection*{RL1.4 Simple Word Boundaries}

\begin{quote}
\emph{To meet this requirement, an implementation shall extend the word
boundary mechanism so that:}

  \begin{enumerate}
  \item
    \emph{The class of \texttt{<word\_character>} includes all the Alphabetic
    values from the Unicode character database, from UnicodeData.txt [UData],
    plus the decimals (General\_Category=Decimal\_Number, or equivalently 
    Numeric\_Type=Deci\-mal), and the U+200C ZERO WIDTH NON-JOINER and
    U+200D ZERO WIDTH JOINER (Join\_Control=True). See also Annex C:
    Compatibility Properties.}

  \item
    \emph{Nonspacing marks are never divided from their base characters, and
    otherwise ignored in locating boundaries.}
  \end{enumerate}
\end{quote}

JFlex does not conform: \verb+\b+ does not match simple word boundaries.

\subsubsection*{RL1.5 Simple Loose Matches}

\begin{quote}
\emph{To meet this requirement, if an implementation provides for 
case-insensitive matching, then it shall provide at least the simple,
default Unicode case-insensitive matching, and specify which properties
are closed and which are not.}
            
\emph{To meet this requirement, if an implementation provides for case
conversions, then it shall provide at least the simple, default Unicode
case folding.}
\end{quote}

JFlex conforms. All supported Unicode Properties are closed.

\subsubsection*{RL1.6 Line Boundaries}

\begin{quote}
\emph{To meet this requirement, if an implementation provides for
line-boundary testing, it shall recognize not only CRLF, LF, CR, but
also NEL (U+0085), PARAGRAPH SEPARATOR (U+2029) and LINE SEPARATOR
(U+2028).}
\end{quote}

JFlex conforms.

\subsubsection*{RL1.7 Supplementary Code Points}

\begin{quote}
\emph{To meet this requirement, an implementation shall handle the full
range of Unicode code points, including values from U+FFFF to U+10FFFF.
In particular, where UTF-16 is used, a sequence consisting of a leading
surrogate followed by a trailing surrogate shall be handled as a single
code point in matching.}
\end{quote}

JFlex does not conform.  Only code points in the Basic Multilingual
Plane (BMP) are supported.  Conformance to RL1.7 is planned for JFlex
\nextver.


\section{A few words on performance}
\label{PerformanceTips}\label{performance}

This section gives some \htmlref{tips}{PerformanceTips} on how to make
your specification produce a faster scanner. 

Although JFlex generated scanners show good performance without
special optimisations, there are some heuristics that can make a
lexical specification produce an even faster scanner. Those are
(roughly in order of performance gain):
 
\begin{itemize}
\item
Avoid rules that require backtracking

From the C/C++ flex \cite{flex} man page: {\trit ``Getting rid
of backtracking is messy and often may be an enormous amount of work for
a complicated scanner.''} Backtracking is introduced by the longest match
rule and occurs for instance on this set of expressions:
 
\noindent
\texttt{\ \ "averylongkeyword"}\\
\texttt{\ \ .}
 
With input \texttt{"averylongjoke"} the scanner has to read all characters
up to \texttt{'j'} to decide that rule \texttt{.} should be matched. All
characters of \texttt{"verylong"} have to be read again for the next
matching process. Backtracking can be avoided in general by adding
error rules that match those error conditions 

\verb+ "av"|"ave"|"avery"|"averyl"|..+

While this is impractical in most scanners, there is still the
possibility to add a ``catch all'' rule for a lengthy list of keywords
\begin{verbatim}
"keyword1"  { return symbol(KEYWORD1); } 
.. 
"keywordn"  { return symbol(KEYWORDn); }
[a-z]+      { error("not a keyword"); }
\end{verbatim}
Most programming language scanners already have a rule like this for
some kind of variable length identifiers.

\item
  Avoid line and column counting
  
  It costs multiple additional comparisons per input character and the
  matched text has to be re-scanned for counting. In most scanners it
  is possible to do the line counting in the specification by
  incrementing \texttt{yyline} each time a line terminator has been
  matched.  Column counting could also be included in actions. This
  will be faster, but can in some cases become quite messy.
 
\item
  Avoid look-ahead expressions and the end of line operator '\$'

  In the best case, the trailing context will first have to be read and 
  then (because it is not to be consumed) re-read again. The cases of 
  fixed-length look-ahead and fixed-length base expressions are handled efficiently
  by matching the concatenation and then pushing back the required amount
  of characters. This extends to the case of a disjunction of fixed-length
  look-ahead expressions such as \verb+r1 / \r|\n|\r\n+. All other cases
  \verb+r1 / r2+ are handled by first scanning the concatenation of 
  \verb+r1+ and \verb+r2+, and then finding the correct end of \verb+r1+. 
  The end of \verb+r1+ is found by scanning forwards in the match again,
  marking all possible \verb+r1+ terminations, and then scanning the reverse
  of \verb+r2+ backwards from the end until a start of \verb+r2+ intersects
  with an end of \verb+r1+. This algorithm is linear in the size of the input
  (not quadratic or worse as backtracking is), but about a factor of 2 slower 
  than normal scanning. It also consumes memory proportional to the size
  of the matched input for \verb+r1 r2+. 

\item
  Avoid the beginning of line operator '\verb+^+'
  
  It costs multiple additional comparisons per match. In some
  cases one extra look-ahead character is needed (when the last character read is
  \verb+\r+, the scanner has to read one character ahead to check if
  the next one is an \verb+\n+ or not).

\item
  Match as much text as possible in a rule.
  
  One rule is matched in the innermost loop of the scanner.  After
  each action some overhead for setting up the internal state of the
   scanner is necessary.
\end{itemize}

Note that writing more rules in a specification does not make the generated
scanner slower (except when you have to switch to another code generation
method because of the larger size).

The two main rules of optimisation apply also for lexical specifications:
\begin{enumerate}
\item {\bf don't do it}
\item {\bf (for experts only) don't do it yet}
\end{enumerate}

Some of the performance tips above contradict a readable and compact
specification style. When in doubt or when requirements are not or not
yet fixed: don't use them --- the specification can always be optimised
in a later state of the development process.


\section{Porting Issues}

\subsection{Porting from JLex\label{Porting}}
JFlex was designed to read old JLex specifications unchanged and to
generate a scanner which behaves exactly the same as the one generated
by JLex with the only difference of being faster.

This works as expected on all well formed JLex specifications.

Since the statement above is somewhat absolute, let's take a look at
what ``well formed'' means here. A JLex specification is well formed, when
it
\begin{itemize}
\item
  generates a working scanner with JLex

\item 
  doesn't contain the unescaped characters \texttt{!} and \texttt{\symbol{126}}

  They are operators in JFlex while JLex treats them as normal
  input characters. You can easily port such a JLex specification
  to JFlex by replacing every \texttt{!} with \verb+\!+ and every
  \verb+~+ with \verb+\~+ in all regular expressions.

\item
  has only complete regular expressions surrounded by parentheses in
  macro definitions
  
  This may sound a bit harsh, but could otherwise be a major problem
  -- it can also help you find some disgusting bugs in your
  specification that didn't show up in the first place. In JLex, a
  right hand side of a macro is just a piece of text, that is copied
  to the point where the macro is used. With this, some weird kind of
  stuff like
  \begin{verbatim}
  macro1 = ("hello"
  macro2 = {macro1})*
  \end{verbatim}  
  was possible (with \texttt{macro2} expanding to \verb+("hello")*+).  This
  is not allowed in JFlex and you will have to transform such
  definitions. There are however some more subtle kinds of errors that
  can be introduced by JLex macros. Let's consider a definition like
  \verb+macro = a|b+  and a usage like \verb+{macro}*+.
  This expands in JLex to \verb+a|b*+ and not to the probably intended
  \verb+(a|b)*+.

  JFlex uses always the second form of expansion, since this is the natural
  form of thinking about abbreviations for regular expressions.

  Most specifications shouldn't suffer from this problem, because
  macros often only contain (harmless) character classes like
  \texttt{alpha = [a-zA-Z]} and more dangerous definitions like

  \verb+ ident = {alpha}({alpha}|{digit})*+

  are only used to write rules like

  \verb+ {ident}       { .. action .. }+

  and not more complex expressions like

  \verb+ {ident}*      { .. action .. }+

  where the kind of error presented above would show up.
\end{itemize}

\subsection{Porting from lex/flex\label{lexport}}
This section tries to give an overview of activities and possible
problems when porting a lexical specification from the C/C++ tools lex
and flex \cite{flex} available on most Unix systems to JFlex.

Most of the C/C++ specific features are naturally not present in JFlex,
but most ``clean'' lex/flex lexical specifications can be ported to 
JFlex without very much work.

This section is by far not complete and is based mainly on a survey of
the flex man page and very little personal experience.  If you do
engage in any porting activity from lex/flex to JFlex and encounter
problems, have better solutions for points presented here or have just
some tips you would like to share, please do 
\latex{contact me via email: \texttt{Gerwin Klein <lsf@jflex.de>}}%
\html{\htmladdnormallink{contact me}{mailto:lsf@jflex.de}}. I will
incorporate your experiences in this manual (with all due credit to you,
of course).

\subsubsection{Basic structure}
A lexical specification for flex has the following basic structure:
\begin{verbatim}
definitions
%%
rules
%%
user code
\end{verbatim}

The \texttt{user code} section usually contains some C code that is used
in actions of the \texttt{rules} part of the specification. For JFlex most
of this code will have to be included in the class code \verb+%{..%}+
directive in the \texttt{options} \texttt{and dec\-la\-ra\-ti\-ons} section (after 
translating the C code to Java, of course). 

\subsubsection{Macros and Regular Expression Syntax}
The \texttt{definitions} section of a flex specification is quite similar
to the \texttt{op\-tions and dec\-la\-ra\-tions} part of JFlex specs.

Macro definitions in flex have the form:
\begin{verbatim}
<identifier>  <expression>
\end{verbatim}
To port them to JFlex macros, just insert a \texttt{=} between \texttt{<identifier>}
and \texttt{<expression>}.

The syntax and semantics of regular expressions in flex are pretty much the
same as in JFlex. A little attention is needed for some escape sequences 
present in flex (such as \verb+\a+) that are not supported in JFlex. These
escape sequences should be transformed into their octal or hexadecimal 
equivalent. 

Another point are predefined character classes. Flex offers the ones directly
supported by C, JFlex offers the ones supported by Java. These classes will
sometimes have to be listed manually (if there is need for this feature, it
may be implemented in a future JFlex version).

\subsubsection{Lexical Rules}
Since flex is mostly Unix based, the '\verb+^+' (beginning of line) and
'\verb+$+' (end of line) operators, consider the \verb+\n+ character as %$
only line terminator. This should usually not cause much problems, but you
should be prepared for occurrences of \verb+\r+ or \verb+\r\n+ or one of
the characters \verb+\u2028+, \verb+\u2029+, \verb+\u000B+, \verb+\u000C+, 
or \verb+\u0085+. They are considered to be line terminators in Unicode and 
therefore may not be consumed when 
\verb+^+ or \verb+$+ is present in a rule.%$

\section{Working together\label{WorkingTog}}

\subsection{JFlex and CUP\label{CUPWork}}
One of the main design goals of JFlex was to make interfacing with the
parser generator CUP \cite{CUP} as easy as possible. 
This has been done by giving
the \texttt{\htmlref{\%cup}{CupMode}} directive a special meaning. However, an
interface always has two sides. This section concentrates on the
CUP side of the story.

\subsubsection{CUP version 0.10j and above}
Since CUP version 0.10j, this has been simplified greatly by the new 
CUP scanner interface \texttt{java\_cup.runtime.Scanner}. JFlex lexers now implement
this interface automatically when the \texttt{\htmlref{\%cup}{CupMode}}
switch is used. There are no special \texttt{parser code}, \texttt{init
  code} or \texttt{scan with} options any more that you have to provide
in your CUP parser specification. You can just concentrate on your grammar.

If your generated lexer has the class name \texttt{Scanner}, the parser
is started from the main program like this:

{\xsmall\begin{verbatim}
...
  try {
    parser p = new parser(new Scanner(new FileReader(fileName)));
    Object result = p.parse().value;
  }
  catch (Exception e) {
...
\end{verbatim}
}

\subsubsection{Custom symbol interface}
If you have used the \texttt{-symbol} command line switch of CUP to change
the name of the generated symbol interface, you have to tell JFlex about
this change of interface so that correct end-of-file code is generated. 
You can do so either by using an \verb+%eofval{+ directive or by using
an \texttt{<<EOF>>} rule. 

If your new symbol interface is called \texttt{mysym} for example, the
corresponding code in the jflex specification would be either

{\xsmall
\begin{verbatim}
%eofval{
  return mysym.EOF;
%eofval}
\end{verbatim}
}

in the macro/directives section of the spec, or it would be

{\xsmall
\begin{verbatim}
  <<EOF>>  { return mysym.EOF; }
\end{verbatim}
}

in the rules section of your spec.

\subsubsection{Using existing JFlex/CUP specifications with CUP 0.10j}
If you already have an existing specification and you would like to upgrade
both JFlex and CUP to their newest version, you will probably have to adjust
your specification.

The main difference between the \texttt{\htmlref{\%cup}{CupMode}} switch in
JFlex 1.2.1 and lower, and the current JFlex version is, that JFlex scanners
now automatically implement the \texttt{java\_cup.runtime.Scanner} interface.
This means the scanning function changes its name from \texttt{yylex()}
to \texttt{next\_token()}. 

The main difference from older CUP versions to 0.10j is, that CUP now
has a default constructor that accepts a \texttt{java\_cup.runtime.Scanner} 
as argument and that uses this scanner as
default (so no \texttt{scan with} code is necessary any more).

If you have an existing CUP specification, it will probably look somewhat like this:
{\xsmall\begin{verbatim}
parser code {:
  Lexer lexer;

  public parser (java.io.Reader input) {
    lexer = new Lexer(input);
  }
:};

scan with {: return lexer.yylex(); :};
\end{verbatim}
}

To upgrade to CUP 0.10j, you could change it to look like this:
{\xsmall\begin{verbatim}
parser code {:
  public parser (java.io.Reader input) {
    super(new Lexer(input));
  }
:};
\end{verbatim}
}

If you do not mind to change the method that is calling the parser,
you could remove the constructor entirely (and if there is nothing else
in it, the whole \texttt{parser code} section as well, of course). The calling
main procedure would then construct the parser as shown in the section above.

The JFlex specification does not need to be changed.


\subsection{JFlex and BYacc/J\label{YaccWork}}

JFlex has built-in support for the Java extension 
\htmladdnormallink{BYacc/J}{http://byaccj.sourceforge.net/} 
\cite{BYaccJ} by Bob Jamison
to the classical Berkeley Yacc parser generator.
This section describes how to interface BYacc/J with JFlex. It
builds on many helpful suggestions and comments from Larry Bell.

Since Yacc's architecture is a bit different from CUP's, the
interface setup also works in a slightly different manner.
BYacc/J expects a function \texttt{int yylex()} in the parser
class that returns each next token. Semantic values are expected
in a field \texttt{yylval} of type \texttt{parserval} where ``\texttt{parser}''
is the name of the generated parser class.

For a small calculator example, one could use a setup like the 
following on the JFlex side:

{\xsmall\begin{verbatim}

%%

%byaccj

%{
  /* store a reference to the parser object */
  private parser yyparser;

  /* constructor taking an additional parser object */
  public Yylex(java.io.Reader r, parser yyparser) {
    this(r);
    this.yyparser = yyparser;
  }
%}

NUM = [0-9]+ ("." [0-9]+)?
NL  = \n | \r | \r\n

%%

/* operators */
"+" | 
..
"(" | 
")"    { return (int) yycharat(0); }

/* newline */
{NL}   { return parser.NL; }

/* float */
{NUM}  { yyparser.yylval = new parserval(Double.parseDouble(yytext()));
         return parser.NUM; }
\end{verbatim}
}

The lexer expects a reference to the parser in its constructor.
Since Yacc allows direct use of terminal characters like \texttt{'+'}
in its specifications, we just return the character code for
single char matches (e.g. the operators in the example). Symbolic
token names are stored as \texttt{public static int} constants in
the generated parser class. They are used as in the \texttt{NL} token
above. Finally, for some tokens, a semantic value may have to be 
communicated to the parser. The \texttt{NUM} rule demonstrates that
bit.

A matching BYacc/J parser specification could look like this:
{\xsmall\begin{verbatim}
%{
  import java.io.*;
%}
      
%token NL          /* newline  */
%token <dval> NUM  /* a number */

%type <dval> exp

%left '-' '+'
..
%right '^'         /* exponentiation */
      
%%

..
      
exp:     NUM          { $$ = $1; }
       | exp '+' exp  { $$ = $1 + $3; }
       ..
       | exp '^' exp  { $$ = Math.pow($1, $3); }
       | '(' exp ')'  { $$ = $2; }
       ;

%%
  /* a reference to the lexer object */
  private Yylex lexer;

  /* interface to the lexer */
  private int yylex () {
    int yyl_return = -1;
    try {
      yyl_return = lexer.yylex();
    }
    catch (IOException e) {
      System.err.println("IO error :"+e);
    }
    return yyl_return;
  }

  /* error reporting */
  public void yyerror (String error) {
    System.err.println ("Error: " + error);
  }

  /* lexer is created in the constructor */
  public parser(Reader r) {
    lexer = new Yylex(r, this);
  }

  /* that's how you use the parser */
  public static void main(String args[]) throws IOException {
    parser yyparser = new parser(new FileReader(args[0]));
    yyparser.yyparse();    
  }
\end{verbatim}
}

Here, the customised part is mostly in the user code section:
We create the lexer in the constructor of the parser and store
a reference to it for later use in the parser's \texttt{int yylex()}
method. This \texttt{yylex} in the parser only calls \texttt{int yylex()}
of the generated lexer and passes the result on. If something goes
wrong, it returns -1 to indicate an error.

Runnable versions of the specifications above 
are located in the \texttt{examples/byaccj} directory of the JFlex 
distribution.

\section{Bugs and Deficiencies\label{Bugs}}

\subsection{Deficiencies}
JFlex \ver conforms with Unicode Regular Expressions
UTS\#18 \cite{unicode_rep} Basic Unicode Support Level 1, with
a few exceptions - for details see 
\htmlref{UTS\#18 Conformance}{unicoderegexconformance}.

\subsection{Bugs}
As of {\today}, no bugs have been reported for JFlex version \ver. All 
bugs reported for earlier versions have been fixed.

If you find new problems, please use the bugs section of the
\htmladdnormallinkfoot{JFlex web site}{http://www.jflex.de/}
to report them.
 
 
\section{Copying and License\label{Copyright}}
JFlex is free software, published under a BSD-style license.

There is absolutely NO WARRANTY for JFlex, its code and its documentation.

See the file \htmladdnormallink{\texttt{COPYRIGHT}}{COPYRIGHT} for more
information.

\newpage  
\begin{thebibliography}{[00]} 
\label{References} 

\bibitem{Aho}
  A.~Aho, R.~Sethi, J.~Ullman, {\trit Compilers: Principles, Techniques, and Tools}, 1986


\bibitem{Appel}
  A.~W.~Appel, {\trit Modern Compiler Implementation in Java: basic techniques}, 1997

\bibitem{JLex}
 E.~Berk, {\trit JLex: A lexical analyser generator for Java},\\
 \htmladdnormallink{\texttt{http://www.cs.princeton.edu/\symbol{126}appel/modern/java/JLex/}}
                   {http://www.cs.princeton.edu/\symbol{126}appel/modern/java/JLex/}


\bibitem{fast}
  K.~Brouwer, W.~Gellerich,E.~Ploedereder, 
  {\trit Myths and Facts about the Efficient Implementation of Finite Automata and Lexical Analysis}, 
  in: Proceedings of the 7th International Conference on Compiler Construction (CC '98), 1998

\bibitem{unicode_rep}
  M.~Davis, {\trit Unicode Regular Expressions}, Unicode Technical Standard \#18, version 17\\ 
  \htmladdnormallink{\texttt{http://www.unicode.org/reports/tr18/tr18-17.html}}
                    {http://www.unicode.org/reports/tr18/tr18-17.html}

\bibitem{ParseTable}
 P.~Dencker, K.~D{\"u}rre, J.~Henft, {\trit Optimization of Parser Tables for portable Compilers}, 
 in: ACM Transactions on Programming Languages and Systems 6(4), 1984

\bibitem{LangSpec}
  J.~Gosling, B.~Joy, G.~Steele, {\trit The Java Language Specifcation}, 1996,\\
  \htmladdnormallink{\texttt{http://java.sun.com/docs/books/jls/}}
                    {http://java.sun.com/docs/books/jls/}

\bibitem{CUP}
  S.~E.~Hudson, {\trit CUP LALR Parser Generator for Java},\\  
  \htmladdnormallink{\texttt{http://www2.cs.tum.edu/projects/cup/}}
  {http://www2.cs.tum.edu/projects/cup/}

\bibitem{BYaccJ}
  B.~Jamison, {\trit BYacc/J},\\
  \htmladdnormallink{\texttt{http://byaccj.sourceforge.net/}}
  {http://byaccj.sourceforge.net}
  
\bibitem{MachineSpec}
 T.~Lindholm, F.~Yellin, {\trit The Java Virtual Machine Specification}, 1996,\\
 \htmladdnormallink{\texttt{http://java.sun.com/docs/books/vmspec/}}
                   {http://java.sun.com/docs/books/vmspec/}
 
\bibitem{flex}
 V.~Paxson, {\trit flex - The fast lexical analyzer generator}, 1995

\bibitem{SparseTable}
  R.~E. Tarjan, A.~Yao, {\trit Storing a Sparse Table}, in: Communications of the ACM 22(11), 1979 

\bibitem{Maurer}
  R.~Wilhelm, D.~Maurer, {\trit {\"U}bersetzerbau}, Berlin $1997^2$

\bibitem{UnicodeSet}
  {\trit Unicode Utilities: UnicodeSet}\\ 
   \htmladdnormallink{\texttt{http://unicode.org/cldr/utility/list-unicodeset.jsp}}
                     {http://unicode.org/cldr/utility/list-unicodeset.jsp}


\end{thebibliography}

\end{document}
